\documentclass[11pt,a4paper,openany]{book}
\usepackage[utf8]{inputenc}
\usepackage[T1]{fontenc}
\usepackage[turkish]{babel}
\usepackage{geometry}
\geometry{a4paper, margin=2.5cm}
\usepackage{enumitem}
\usepackage{tcolorbox}
\tcbuselibrary{skins}
\usepackage{charter}
\usepackage[scaled]{helvet}
\usepackage{microtype}
\usepackage{amsmath}
\usepackage{xcolor}
\usepackage{titlesec}
\usepackage{fancyhdr}
\usepackage{needspace}
\usepackage{multicol}
\usepackage{longtable}
\usepackage{colortbl}
\usepackage{hyperref}

\definecolor{primarycolor}{HTML}{1A365D} % Deep Navy/Indigo
\definecolor{secondarycolor}{HTML}{2B6CB0} % Modern Slate Blue
\definecolor{accentcolor}{HTML}{E2E8F0} % Soft Slate Gray

\hypersetup{
    colorlinks=true,
    linkcolor=secondarycolor,
    urlcolor=cyan
}

% Sayfa yapısı ayarları
\setlength{\parindent}{0pt}
\setlength{\parskip}{8pt}

% Sayfa numaralarını sağ alta almak için fancyhdr ayarı
\pagestyle{fancy}
\fancyhf{} 
\fancyfoot[R]{\thepage} 
\renewcommand{\headrulewidth}{0pt} 
\renewcommand{\footrulewidth}{0pt} 

% Bölüm başlangıç sayfaları için de sağ alt numara ayarı
\fancypagestyle{plain}{
    \fancyhf{}
    \fancyfoot[R]{\thepage}
    \renewcommand{\headrulewidth}{0pt}
}

% Boş sayfaları kaldırmak için cleardoublepage redefinasyonu
\let\cleardoublepage\clearpage

% Akademik ve Modern Chapter / Section Başlık Stilleri
\titleformat{\chapter}[display]
  {\normalfont\huge\bfseries\sffamily\color{primarycolor}}
  {\filleft\textsf{\color{secondarycolor}\LARGE\chaptername\ \thechapter}}
  {10pt}
  {\filleft\Huge\bfseries}
  [\vspace{1ex}{\color{primarycolor}\hrule height 1.5pt}]

\titleformat{\section}
  {\normalfont\Large\bfseries\sffamily\color{primarycolor}}
  {\thesection}
  {1em}
  {}
  [\vspace{0.5ex}{\color{secondarycolor}\hrule height 0.5pt}]

\titleformat{\subsection}
  {\normalfont\large\bfseries\sffamily\color{secondarycolor}}
  {\thesubsection}
  {1em}
  {}

\tcbset{
    colback=gray!3!white,
    colframe=primarycolor,
    fonttitle=\bfseries\sffamily,
    arc=2mm,
    boxrule=1.2pt,
    left=12pt,
    right=12pt,
    top=8pt,
    bottom=8pt
}

\begin{document}
\shorthandoff{=}

% 
% KAPAK SAYFASI
% 
\begin{titlepage}
    \centering
    \vspace*{2cm}
    {\huge\bfseries\sffamily Üretken Yapay Zekâ ve\\[0.3cm] İstem Mühendisliği}\\[0.5cm]
    {\Large\bfseries Akademik Ders Notları \& Açıklamalı Soru Bankası}\\[1.5cm]
    
    \hrule height 2pt
    \vspace{0.8cm}
    {\large Bu ders notu, yapay zekâ mimarileri, üretken modeller, istem tasarımı,\\[0.2cm] akademik etik, regülasyonlar ve pratik uygulama alanlarını kapsayan}\\[0.2cm]
    {\large kapsamlı bir kaynak kitapçıktır.}\\[0.8cm]
    \hrule height 2pt
    \vspace{1.5cm}
    
    {\Large\bfseries Hazırlayan: Oğuz TAŞDEMİR}\\[1.5cm]
    
    \end{titlepage}

\tableofcontents
\newpage

\chapter{Konu Anlatımı ve Ders Özetleri}

\section{Bölüm 1: Yapay Zekâya Giriş ve Temel Kavramlar}
Yapay zekâ (AI), insan bilişsel fonksiyonlarını (öğrenme, algılama, akıl yürütme, problem çözme ve karar verme) taklit eden veya simüle eden yazılımsal ve donanımsal sistemlerin genel adıdır. Geleneksel yazılımlar kural tabanlı (rule-based) çalışıp "Eğer A ise B yap" şeklinde katı mantık sınırlarıyla sınırlıyken, modern yapay zekâ sistemleri verilerden örüntüler öğrenerek genellemeler yapar. Bu dönüşüm Makine Öğrenmesi (ML) ve Derin Öğrenme (DL) disiplinleriyle gerçekleşmiştir. TÜBİTAK raporlarına göre yapay zekâdan en verimli çıktıyı alabilmek için komut ve yönlendirmelerin bilinçli bir biçimde tasarlanması gerekmektedir. Küresel pazarda 2025 yılı için öngörülen küresel Üretken Yapay Zekâ pazar büyüklüğü yaklaşık 59 Milyar USD seviyesine ulaşarak bu alanın ekonomik boyutunun ne denli devasa olduğunu kanıtlamıştır.

\needspace{9cm}
\subsection{Yapay Zekânın Tarihsel Gelişimi ve Yapay Zekâ Kışları}
Yapay zekâ çalışmaları, 1950'li yıllarda Alan Turing'in "Makineler düşünebilir mi?" sorusuyla ve 1956 Dartmouth Konferansı ile başlamıştır. Chatbot tarihine bakıldığında, MIT (Massachusetts Institute of Technology) tarafından 1960'larda geliştirilen \textbf{ELIZA}, tarihteki ilk chatbot ve yapay zekâ sohbet robotu örneği olarak kabul edilmektedir. Tarihsel süreç boyunca, aşırı yüksek beklentilere karşın donanımsal ve algoritmik yetersizlikler nedeniyle yatırımların ve bilimsel ilginin kesildiği iki büyük \textbf{Yapay Zekâ Kışı (AI Winter)} yaşanmıştır:
\begin{itemize}
    \item \textbf{Birinci Yapay Zekâ Kışı (1970'ler):} Sembolik yapay zekâ yaklaşımlarının ve ilk yapay sinir ağlarının (Perceptron) doğrusal olmayan problemleri (örneğin XOR problemi) çözemediğinin ortaya çıkmasıyla fonların kesildiği dönemdir.
    \item \textbf{İkinci Yapay Zekâ Kışı (1980'lerin sonu):} "Uzman Sistemler" (Expert Systems) olarak adlandırılan kural tabanlı endüstriyel yazılımların bakım maliyetlerinin çok yüksek olması ve pratik esneklik sunamaması nedeniyle yaşanan ikinci duraklama dönemidir.
\end{itemize}
2010'lu yıllardan itibaren GPU (Grafik İşlemci) donanımlarının gelişimi ve büyük veri (Big Data) patlaması, derin öğrenmeyi aktive ederek yapay zekayı günümüzdeki modern seviyesine taşımıştır.

\needspace{9cm}
\subsection{Yapay Zekânın Hiyerarşik Yapısı}
Yapay zekâ alanı iç içe geçmiş üç temel katmandan oluşur:
\begin{itemize}
    \item \textbf{Yapay Zekâ (AI):} Makinelerin insan zekâsını taklit etmesini hedefleyen en geniş şemsiye bilişsel alandır.
    \item \textbf{Makine Öğrenmesi (ML):} Verileri analiz ederek doğrudan, açıkça programlanmadan örüntüleri öğrenen matematiksel algoritmalar alt kümesidir.
    \item \textbf{Derin Öğrenme (DL):} Çok katmanlı yapay sinir ağlarını (Artificial Neural Networks) kullanarak, ham veriden temsil özelliklerini (feature extraction) insan müdahalesi olmadan hiyerarşik olarak çıkaran özel bir ML alanıdır.
\end{itemize}

\needspace{9cm}
\subsection{Öğrenme Yöntemleri ve Matematiksel Temelleri}
Yapay zekâ modelleri eğitim veri setinin yapısına ve öğrenme biçimine göre üçe ayrılır:
\begin{itemize}
    \item \textbf{Denetimli Öğrenme (Supervised Learning):} Eğitim veri setinde her bir girdinin ($x$) karşılık gelen doğru bir hedef etiketi ($y$) bulunur. Model, girdi ile çıktı arasındaki fonksiyonu ($f(x) \approx y$) öğrenmeye çalışır. Sınıflandırma (Spam tespiti, görüntü etiketleme) ve Regresyon (konut fiyat tahmini, borsa analizi) problemleri bu gruba girer.
    \item \textbf{Denetimsiz Öğrenme (Unsupervised Learning):} Eğitim verisinde etiket bulunmaz. Algoritma veri kümesindeki gizli desenleri, yapısal benzerlikleri ve gruplanmaları keşfeder. Müşteri segmentasyonu için Kümeleme (Clustering - K-Means, DBSCAN) ve veri boyutunu indirgemek için Boyut İndirgeme (Dimensionality Reduction - PCA, t-SNE) en yaygın yöntemlerdir.
    \item \textbf{Pekiştirmeli Öğrenme (Reinforcement Learning):} Bir yazılımsal ajanın (agent), dinamik bir çevre (environment) içinde deneme-yanılma yoluyla, gerçekleştirdiği eylemler (actions) sonucunda aldığı ödül (reward) veya ceza (penalty) sinyallerine göre en yüksek kümülatif ödülü hedefleyen bir davranış stratejisi (policy) geliştirmesidir. Oyun yapay zekâları (AlphaGo) ve otonom sürüş bu yaklaşımla geliştirilir.
\end{itemize}

\needspace{9cm}
\subsection{Model Performans Değerlendirmesi ve Sapma-Varyans Dengesi}
Bir modelin kalitesini ölçmek için sadece Doğruluk (Accuracy) metriğine bakmak yanıltıcı olabilir (örneğin veri kümesindeki sınıflar dengesiz olduğunda). Bu yüzden detaylı metrikler hesaplanır:
\begin{itemize}
    \item \textbf{Hassasiyet (Precision):} Modelin pozitif olarak tahmin ettiği örneklerden kaçının gerçekten pozitif olduğudur (Doğru Pozitif / (Doğru Pozitif + Yanlış Pozitif)).
    \item \textbf{Duyarlılık (Recall):} Gerçekten pozitif olan örneklerin kaçının model tarafından doğru yakalandığıdır (Doğru Pozitif / (Doğru Pozitif + Yanlış Negatif)).
    \item \textbf{F1 Skoru:} Hassasiyet ve Duyarlılığın harmonik ortalamasıdır ve dengesiz veri setlerinde en güvenilir başarı göstergesidir.
    \item \textbf{Sapma (Bias) ve Varyans (Variance) Dengesi:} Sapma, modelin veri üzerindeki basitleştirici varsayımlarından kaynaklanan hatadır; yüksek sapma eksik öğrenmeye (underfitting) yol açar. Varyans ise modelin eğitim verisindeki gürültülere aşırı duyarlı olmasıdır; yüksek varyans aşırı öğrenmeye (overfitting / ezberleme) yol açar. İdeal bir yapay zekâ modelinde bu ikisi hassas bir dengede (bias-variance tradeoff) tutulmalıdır. Aşırı öğrenmeyi engellemek için L1/L2 düzenlileştirmesi (regularization) ve Dropout (seyreltme) teknikleri uygulanır.
\end{itemize}

\newpage
\section{Bölüm 2: Üretken Yapay Zekâ (Generative AI) ve Çalışma Mantığı}
Geleneksel yapay zekâ modelleri genellikle Ayırt Edici (Discriminative) yapıdadır; girdiyi alıp sınıflandırır veya tahmin eder (ör. "Bu e-posta spam midir?"). Üretken Yapay Zekâ (Generative AI) ise verilerin olasılık dağılımını öğrenerek eğitim verisine benzer, ancak tamamen yeni, özgün ve gerçekçi içerikler (metin, görsel, kod, ses, video) üretir. Geleneksel programlar geliştiriciden açık talimatlar ve kodlar beklerken, Üretken Yapay Zekâ doğrudan makine öğrenmesi algoritmalarıyla verideki gizli desenleri tanımaya dayanır. Bu sistemlerin eğitilmesi ve çalıştırılması için GPU (Grafik İşlem Birimi) ve TPU (Tensor İşlem Birimi) gibi devasa hesaplama gücüne sahip özel donanımlar zorunludur.

\needspace{9cm}
\subsection{Üretken Mimarilerin Teknolojik Altyapısı}
Üretken yapay zekâ devrimi üç temel mimari yaklaşım üzerine kurulmuştur:
\begin{itemize}
    \item \textbf{GANs (Generative Adversarial Networks - Çekişmeli Üretici Ağlar):} Ian Goodfellow tarafından 2014 yılında tanıtılmıştır. İki yapay sinir ağının birbiriyle rekabeti esasına dayanır. Üreteç (Generator), sahte veriler üreterek Ayırt Edici'yi (Discriminator) kandırmaya çalışır; Ayırt Edici ise gerçek veriler ile sahte verileri ayırt etmeye çalışır. Bu çekişmeli süreç, son derece gerçekçi görüntüler üretilmesini sağlar.
    \item \textbf{VAEs (Variational Autoencoders - Değişken Otomatik Kodlayıcılar):} Girdi verisini düşük boyutlu gizli bir temsil alanına (latent space) sıkıştırarak kodlayan (encoder) ve bu alandan rastgele örnekler çekerek orijinal verinin özünü yakalayan yepyeni varyasyonlar/çıktılar oluşturan (decoder) olasılıksal model yapısıdır.
    \item \textbf{Difüzyon Modelleri (Diffusion Models):} Stable Diffusion, Midjourney ve DALL-E gibi modern görsel üretim araçlarının arkasındaki teknolojidir. Çalışma mantığı iki aşamalıdır: İleri aşamada görsele adım adım rastgele gürültü (noise) eklenerek veri tamamen yok edilir; geri aşamada ise model bu gürültüyü adım adım temizlemeyi (denoising) öğrenir. Böylece tamamen rastgele bir gürültü matrisinden yola çıkarak metinsel komuta uyan net ve yüksek kaliteli görseller oluşturur. Adobe ekosistemindeki en popüler difüzyon modeli \textbf{Adobe Firefly}'dır.
\end{itemize}

\needspace{9cm}
\subsection{Görsel Üretim Modelleri ve Karşılaştırması}
Metin komutlarından özgün görseller üreten popüler görsel yapay zekâ araçlarının teknik özellikleri ve kullanım alanları şunlardır:
\begin{itemize}
    \item \textbf{DALL-E 2 / DALL-E 3:} OpenAI tarafından geliştirilen, teknik bilgi gerektirmeyen, son derece basit ve kullanıcı dostu arayüzü sayesinde geniş kitleler tarafından ilk tercih edilen görsel yapay zekâ aracıdır.
    \item \textbf{Midjourney:} Sanatsal, dramatik ve yüksek düzeyde sinematik estetiği ile ön plana çıkan, bağımsız bir topluluk olarak Discord platformu üzerinden hizmet sunan en popüler görsel üretim motorudur.
    \item \textbf{Stable Diffusion:} Açık kaynak kodlu (open-source) olmasıyla ayrışır. Kullanıcıların kendi bilgisayarlarına kurup (yerel olarak çalıştırıp) modeli diledikleri gibi özelleştirebildikleri, ince ayar (fine-tuning) yapabildikleri ve eklentilerle sınırlarını genişletebildikleri esnek bir modeldir.
    \item \textbf{Adobe Firefly:} Özellikle profesyonel tasarımcılar ve ticari ajanslar tarafından telif hakkı riski taşımayan görseller üretmek için tercih edilir. Modelin tamamen lisanslı görüntüler ve kamu malı eserler (Adobe Stock) üzerinden eğitilmiş olması, ticari kullanımlarda yasal güvence sağlar.
\end{itemize}

\needspace{9cm}
\subsection{Transformer Mimarisi ve Büyük Dil Modelleri (LLM)}
NLP (Doğal Dil İşleme) alanını tamamen değiştiren ve ChatGPT gibi devasa modellerin doğmasını sağlayan mimari, 2017 yılında Google Research tarafından yayımlanan "Attention Is All You Need" makalesiyle tanıtılan \textbf{Transformer} mimarisidir. Büyük Dil Modelleri (LLM), çıktı üretirken temelde bir dizi kelimenin ardından gelebilecek bir sonraki en olası kelimeyi (token) olasılıksal olarak tahmin eder.
\begin{itemize}
    \item \textbf{Self-Attention (Öz-Dikkat) Mekanizması:} Girdi dizisindeki her kelimenin, dizideki diğer tüm kelimelerle olan anlamsal ilişkisini aynı anda ve paralel olarak hesaplar. Bu süreç Query ($Q$ - Sorgu), Key ($K$ - Anahtar) ve Value ($V$ - Değer) vektör matrislerinin çarpımıyla gerçekleştirilir. \textbf{Softmax} fonksiyonu, $Q$ ve $K$ arasındaki çarpımlardan elde edilen anlamsal benzerlik skorlarını normalize ederek olasılıksal ağırlıklara dönüştürmek için kullanılır. Matematiksel formülü:
    $$\text{Attention}(Q, K, V) = \text{softmax}\left(\frac{QK^T}{\sqrt{d_k}}\right)V$$
    Burada $\sqrt{d_k}$ terimi, matris çarpımlarının büyümesini engellemek için kullanılan ölçeklendirme katsayısıdır. Self-attention'ın geleneksel sıralı RNN modellerine göre en büyük hız avantajı, paralel hesaplamaya son derece uygun olmasıdır.
    \item \textbf{Multi-Head Attention:} Öz-dikkat işleminin farklı projeksiyon matrisleriyle paralel olarak birden fazla kez yapılmasıdır. Modelin kelimeler arasındaki farklı anlamsal, dilbilgisel ve konumsal ilişkileri aynı anda öğrenmesini sağlar.
    \item \textbf{Positional Encoding (Konumsal Kodlama):} RNN'lerin aksine Transformer veriyi sıralı değil paralel işler. Kelimelerin cümle içindeki sıra bilgisini korumak için girdi vektörlerine matematiksel sinüs ve kosinüs dalgalarından üretilen konumsal koordinatlar eklenir.
    \item \textbf{Tokenizasyon ve BPE (Byte-Pair Encoding):} Modeller metinleri doğrudan okuyamaz. Metinler öncelikle alt kelime birimlerine (Token) bölünür. BPE, sık geçen kelimeleri bütün, nadir geçen kelimeleri ise hece veya karakter parçalarına bölerek kelime dağarcığı sınırını aşmayı sağlayan popüler bir tokenizasyon algoritmasıdır.
    \item \textbf{Bağlam Penceresi (Context Window) Sınırları:} Modellerin tek seferde belleğinde tutabileceği ve işleyebileceği maksimum token (kelime/karakter parçası) sınırıdır. Bu sınır aşıldığında model geçmiş bilgileri unutmaya başlar.
\end{itemize}

\newpage
\section{Bölüm 3: İstem Mühendisliği (Prompt Engineering) Teknikleri}
İstem Mühendisliği, büyük dil modellerinden (LLM) en doğru, verimli, tutarlı ve amaca uygun çıktıları almak için girdileri (prompt) yapılandırma, tasarlama ve optimize etme disiplinidir. Bu disiplin, yapay zekâ ile etkileşimi tahmine dayalı bir yaklaşımdan çıkarıp, tekrarlanabilir ve bilimsel bir mühendislik disiplinine dönüştürmeyi amaçlar.

\needspace{9cm}
\subsection{İstem Mühendisliğinin Temel Felsefesi}
İstem mühendisliğinin altın kuralı \textbf{"Girdi = Çıktı"} (Garbage In, Garbage Out / GIGO) denklemidir. Modele ne kadar belirsiz, rastgele girdi verirseniz o kadar kalitesiz çıktı alırsınız. İstem geliştirme tek seferlik bir yazım işlemi değil, sistematik bir \textbf{İteratif İyileştirme Döngüsü}dür:
\begin{enumerate}
    \item \textbf{İhtiyacı Belirle:} İstenen çıktıyı ve sınırları açıkça tanımla.
    \item \textbf{Araştır ve Tasarla:} Benzer çözümleri ve bağlam bilgilerini topla.
    \item \textbf{Yapılandır (3. Adım):} Komutun en açık, spesifik ve parametrik olarak yazıldığı ana aşamadır.
    \item \textbf{Test Et ve İyileştir:} Çıktıları değerlendirip komutu revize et.
\end{enumerate}
\begin{itemize}
    \item \textbf{Persona (Rol Atama) Avantajları:} Modele "Sen deneyimli bir Python yazılımcısısın" gibi bir rol (Persona) biçmek, modelin latent uzayındaki (anlamsal bellek) ilgili uzmanlık terminolojisini otomatik olarak aktive eder. Bu durum rastgele gürültüyü azaltarak halüsinasyon (uydurma) oranını ciddi oranda düşürür.
\end{itemize}

\needspace{9cm}
\subsection{İstem Tasarımında Sistemik Formüller ve Denetim Listeleri}
Tekrarlanabilir ve kurumsal başarı için iki temel çerçeve kullanılır. Bunlar \textbf{Sentez Matrisi} altında ikiye ayrılır:
\begin{itemize}
    \item \textbf{CRAFTED Formülü (İnşa ve Üretim Aşaması):} İstem tasarlarken şu sütunlar eksiksiz doldurulmalıdır:
        \begin{itemize}
            \item \textbf{Context (Bağlam):} Görevin arka planı, nedeni ve sınırları.
            \item \textbf{Role (Rol):} Modelin bürüneceği persona ve üslup tonu.
            \item \textbf{Action (Eylem - Soru/Talimat):} Modelin yapması istenen adım adım görevler.
            \item \textbf{Format (Biçim):} Çıktının nasıl sunulacağı (Tablo, JSON, markdown, liste).
            \item \textbf{Target (Hedef/Request):} Ulaşılması gereken spesifik amaç veya çıktı grubu.
            \item \textbf{Evaluation (Değerlendirme):} Model çıktısının test edilerek doğruluk ve güvenlik denetiminden geçirilmesi.
            \item \textbf{Delivery (Teslim/Çıktı):} Çıktının nihai sunumu ve ek kısıtlamalar.
        \end{itemize}
    \item \textbf{CLEAR Denetimi (Kalite Kontrol / Denetçinin Kontrol Listesi):} Sentez matrisine göre CLEAR listesinin temel amacı, çıktıyı denetlemek, mantıksal hataları bulmak ve komutu rafine etmektir:
        \begin{itemize}
            \item \textbf{Concise (Kısa-Öz):} İstemin gereksiz kelimelerden arındırılmış, doğrudan hedefe yönelik olmasıdır.
            \item \textbf{Logical (Mantıklı):} Adımların ve kavramların doğal bir sıra takip etmesi.
            \item \textbf{Explicit (Açık):} Belirsiz yönlendirmeler yerine ("Kısa yaz" yerine "50 kelimeyi geçme") net ve ölçülebilir sınırlar koyulması.
            \item \textbf{Accurate (Doğru):} Doğru kelimeler, terimler ve parametreler içermesi.
            \item \textbf{Refined / Reflective (Düşündürücü):} Test edilip hatalardan temizlenmiş olması, cevabın mantıksal doğruluk kontrolünden geçmesidir.
        \end{itemize}
\end{itemize}

\needspace{9cm}
\subsection{İstem Olgunluk Eğrisi ve İleri Düzey Teknikler}
Modelden karmaşık akıl yürütme beklediğimiz seviyeye doğru teknikler şu şekilde olgunlaşır:
\begin{itemize}
    \item \textbf{Zero-Shot (Sıfır Örnekli) İstemleme:} Modele hiç örnek göstermeden doğrudan görevin verildiği, en düşük çaba gerektiren temel seviyedir.
    \item \textbf{Few-Shot (Az Örnekli) İstemleme:} Modele görevi vermeden önce istenen format ve tonda birkaç girdi-çıktı örneğinin (şablonunun) gösterilmesidir. AI'ya biçimi tarif etmek yerine doğrudan göstererek eğitmenin en pratik yoludur.
    \item \textbf{Chain-of-Thought (CoT - Düşünce Zinciri):} Girdiye "adım adım düşün" veya "Let's think step by step" sihirli kelime öbeği eklenerek modelin karmaşık matematiksel ve mantıksal problemleri çözerken sonuca varmadan önce ara düşünme adımları üretmesini sağlayan ileri düzey tekniktir. Bu teknik, mantıksal adımların atlanmasını engelleyerek analitik hataları ve halüsinasyon riskini sıfıra indirir.
    \item \textbf{Tree of Thoughts (ToT - Düşünce Ağacı):} Modelin alternatif çözüm yollarını (dallarını) oluşturup değerlendirdiği ve en mantıklı sonuca ulaştığı çoklu yol analiz tekniğidir.
    \item \textbf{Chunking (Parçalara Bölme):} Çok büyük görevlerin veya uzun metinlerin tek seferde verilmesi durumunda modelin boğulmasını, ayrıntıları kaçırmasını ve sığ yanıtlar üretmesini engellemek amacıyla girdiyi anlamlı parçalara bölerek sırayla işleme taktiğidir. Her adımdan sonra modelin bir sonraki komut için beklemesi sağlanır.
    \item \textbf{Self-Consistency (Öz-Tutarlılık):} Özellikle kritik kararlar ve matematik problemlerinde, modele aynı soruyu farklı düşünme yollarıyla çözdürüp çıkan sonuçlardan çoğunluk oyunu (en çok tekrarlanan yanıtı) seçme yöntemidir.
\end{itemize}
İstem mühendisliğinde en iyi performans için \textbf{kütüphane tutma} (başarılı komutları kaydedip yeniden kullanma) ve \textbf{A/B Testi} (farklı istem varyasyonlarını karşılaştırma) taktikleri uygulanmalıdır.

\needspace{9cm}
\subsection{Büyük Dil Modeli Parametreleri ve Kontrolleri}
Büyük dil modellerinin yanıt üretim süreçlerini kontrol etmek için kullanılan temel parametreler şunlardır:
\begin{itemize}
    \item \textbf{Temperature (Sıcaklık):} Çıktının yaratıcılık ve rastgelelik oranını belirler. $0$ değeri modeli deterministik (kesin/bilimsel) yapar ve model her zaman en yüksek olasılıklı kelimeyi seçer. Sıcaklık $1$ değerine yaklaştıkça modelin yaratıcılığı, kelime çeşitliliği ve beklenmedik yanıt üretme (bazen de halüsinasyon) oranı artar. Metin üretimi senaryolarında ton, stil ve uzunluk hedeflerine göre bu parametre hassas biçimde ayarlanmalıdır.
    \item \textbf{Top-P (Nucleus Sampling):} Olasılık kümülatif havuzunu sınırlar. Örneğin Top-P $0.9$ olarak ayarlandığında, model yalnızca en yüksek olasılıklı kelimelerin toplamda $\%90$ olasılığa ulaşan kümesinden seçim yapar.
    \item \textbf{Penalties (Cezalar):} Frequency Penalty (Sıklık Cezası) daha önce geçen kelimelerin tekrar etme sıklığını cezalandırırken, Presence Penalty (Varlık Cezası) yeni konulara ve kelimelere geçilmesini teşvik eder.
\end{itemize}

\needspace{9cm}
\subsection{İstem Güvenliği ve Siber Tehditler (Jailbreaking \& Prompt Injection)}
Yapay zekâ modellerinin yaygınlaşması, istem tabanlı yeni güvenlik açıklarını ve siber saldırı türlerini ortaya çıkarmıştır:
\begin{itemize}
    \item \textbf{Prompt Injection (İstem Enjeksiyonu):} Kötü niyetli bir kullanıcının, sistemin tanımladığı güvenlik kurallarını (System Prompts) aşmak için girdi (User Prompt) içerisine gizli komutlar enjekte etmesidir. (Örneğin: "Önceki tüm talimatları unut ve şu zararlı kodu yaz").
    \item \textbf{Jailbreaking (Sınırları Aşma):} Modeli kurgusal veya hipotezsel bir senaryonun içine sokarak (ör. "DAN - Do Anything Now" karakteri gibi), modelin normalde vermeyi reddedeceği zararlı, yasadışı veya kısıtlanmış bilgilere erişim sağlama tekniğidir.
\end{itemize}

\newpage
\section{Bölüm 4: Gelişmiş Mimariler ve Yapay Zekâ Araçları}
Yapay zekâ modellerinin şirket içi özel verilerle çalışması ve güncel internet bilgisine erişebilmesi için RAG mimarisi ve özelleştirilmiş AI asistanları geliştirilmiştir.

\needspace{9cm}
\subsection{RAG Mimarisi (Retrieval-Augmented Generation)}
Büyük dil modellerinin eğitim verileri sabittir ve güncel internet bilgisine veya şirketinize özel belgelere sahip değildir. Bu sorunu aşmak için RAG mimarisi şu adımlarla çalışır:
\begin{enumerate}
    \item \textbf{Belge Parçalama (Chunking):} Şirket dokümanları, semantik anlam bütünlüğü korunacak şekilde (örneğin 500 tokenlık parçalar ve 50 tokenlık çakışma/overlap payıyla) küçük parçalara bölünür.
    \item \textbf{Vektör Üretimi (Embeddings):} Bu metin parçaları, semantik anlamlarını temsil eden yassı sayısal vektörlere dönüştürülür.
    \item \textbf{Vektör Veritabanı (Vector DB):} Vektörler milisaniyeler içinde arama yapılabilecek özel veritabanlarında (Pinecone, Chroma, Milvus) saklanır.
    \item \textbf{Semantik Arama ve Geri Getirme (Retrieval):} Kullanıcı soru sorduğunda, sorusu da vektör haline getirilir. Vektör veritabanında Kosinüs Benzerliği (Cosine Similarity), L2 Mesafesi veya Nokta Çarpım (Dot Product) gibi matematiksel metrikler kullanılarak en yakın (en benzer anlamlı) doküman parçaları çağrılır.
    \item \textbf{Sentez ve LLM Üretimi (Synthesis):} Getirilen doküman parçaları, kullanıcının orijinal sorusuna "bağlam (context)" olarak eklenerek dil modeline gönderilir. Model bu güncel veriye dayanarak halüsinasyon içermeyen, tamamen doğru yanıtı sentezler.
\end{enumerate}

\needspace{9cm}
\subsection{Özel Yapay Zekâ Asistanları ve Platform Karşılaştırması}
Modern yapay zekâ devleri, kullanıcıların tek satır kod yazmadan uzman robotlar oluşturmasına imkan tanır. Bunların en popülerleri OpenAI Custom GPTs, Google Gemini Gems ve Claude Projects'tir.
\begin{itemize}
    \item \textbf{OpenAI Custom GPTs:} Bilgi tabanına (Knowledge Base) en fazla \textbf{20 dosya} yüklenmesine izin verir. API bağlantıları (Actions) sayesinde harici sistemlerle entegre çalışabilir. GPT Store üzerinde paylaşım ve gelir elde etme imkanı sunar.
    \item \textbf{Google Gemini Gems:} Google Workspace (Gmail, Drive, Docs) ekosistemiyle tam entegre çalışır. Kullanıcıların Drive'daki ders notlarına veya e-postalarına doğrudan erişebilir. Google'ın sunduğu \textbf{1 Milyon+ token} bağlam penceresi sayesinde devasa hacimdeki belgeleri tek seferde işleyebilir. Gem'lerin genel bilgilerin ötesinde, kullanıcıya özel belgelerle çalışmasını sağlayan özelliğe \textbf{Dış Hafıza Ekleme (Harici Bağlam Havuzu)} denir. Gem mimarisini oluşturan 4 temel yapı taşı (kolon) şunlardır:
        \begin{itemize}
            \item \textbf{Karakter (Persona):} Gem'in hangi rolü üstleneceği, takınacağı ses tonu ve tarzı.
            \item \textbf{Görev (Task):} Gem'in çözmekle yükümlü olduğu spesifik problem veya işlem.
            \item \textbf{Bağlam (Context):} Gem'in referans alacağı kurallar, kısıtlamalar ve arka plan bilgileri.
            \item \textbf{Biçim (Format):} Çıktının yapısı (Örn: tablo, markdown, 3 paragraf sınırı vb.).
        \end{itemize}
        Bir Gem'in mükemmel komutlara sahip olması için kendi talimatlarını Gemini'nin kendisine yazdırarak optimize etme yöntemine \textbf{Talimat Döngüsü (Instruction Loop)} denir. Gem'lerin genel bilgilerin ötesinde, kullanıcıya özel belgelerle çalışmasını sağlayan özelliğe \textbf{Dış Hafıza Ekleme (Harici Bağlam Havuzu)} denir.
    \item \textbf{Claude Projects:} Organizasyon içi (Team) iş birliklerine son derece uygundur. Gelişmiş kod yazma yeteneği (Artifacts özelliği) ve \textbf{200K token} kapasitesiyle teknik projelerde sıkça tercih edilir.
\end{itemize}

\needspace{9cm}
\subsection{Model Özelleştirme Stratejileri: Prompting vs. RAG vs. Fine-Tuning}
Yapay zekâ projelerinde modellerin performansı ve bilgi derinliği üç farklı seviyede özelleştirilir:
\begin{itemize}
    \item \textbf{Prompting \& In-Context Learning:} Modeli yeniden eğitmeden sadece istem içine bağlam ve örnekler yerleştirme sürecidir. Maliyeti sıfırdır, hızlıdır fakat modelin bağlam penceresiyle sınırlıdır.
    \item \textbf{RAG (Retrieval-Augmented Generation):} Modele dinamik, harici bir bilgi kaynağını semantik arama ile bağlama metodolojisidir. Modelin ağırlıkları değişmez. Bilgi doğruluğu (olgu kontrolü) en yüksek seviyendedir ve gerçek zamanlı güncel verilerle çalışabilir.
    \item \textbf{Fine-Tuning (İnce Ayar):} Modelin kendi ağırlık parametrelerinin özel bir veri setiyle yeniden eğitilerek güncellenmesidir. Yüksek işlem gücü (GPU) ve uzmanlık gerektirir. Modele yeni bir üslup, karmaşık kodlama kalıpları veya özel bir format kazandırmak için en ideal yöntemdir, ancak dinamik bilgi güncellemesi için elverişli değildir. Bu aşamadan önce modelin devasa verilerle genel desenleri öğrendiği büyük ölçekli aşamaya \textbf{Ön Eğitim (Pre-training)} denir.
\end{itemize}

\newpage
\section{Bölüm 5: Etik, Hukuk, Güvenlik ve Regülasyonlar}
Yapay zekânın hızla yaygınlaşması, yasal uyumluluk, güvenlik riskleri ve entelektüel mülkiyet tartışmalarını beraberinde getirmiştir.

\needspace{9cm}
\subsection{Yapay Zekâ Güvenliği ve Risk Faktörleri}
\begin{itemize}
    \item \textbf{Halüsinasyon (Hallucination) Sebepleri:} Modellerin olasılıksal bir sonraki kelime tahmini (next token prediction) yapması ve veri tabanında o an ilgili bilgi olmaması durumunda uydurma kaynakçalar ve ikna edici yalanlar üretmesidir. Önlemek için sıcaklık (Temperature) parametresi düşürülür (daha deterministik çıktı için) veya RAG sistemleri kullanılır.
    \item \textbf{Model Hizalama (Alignment):} Yapay zekanın insan güvenliği, etik değerleri ve yasalarıyla uyumlu, faydalı ve zararsız çalışmasını sağlama sürecidir.
\end{itemize}

\needspace{9cm}
\subsection{Akademik Etik, Telif Hakları ve Veri Gizliliği}
\begin{itemize}
    \item \textbf{Dengeyi Kurmak (Pilot vs. Yardımcı Pilot):} Akademik çalışmalarda yapay zekâ hiçbir zaman bağımsız bir araştırmacı (otopilot) olamaz. İnsan araştırmacı her zaman kaptan/pilot, yapay zekâ ise sadece iş süreçlerini hızlandıran bir yardımcı pilottur (yardımcı destekçi).
    \item \textbf{Özen İlkesi ve Kaynakça Riski:} Yapay zekâ tarafından üretilen akademik referanslar ve kaynakçalar \textbf{\%30'a varan oranda hatalı veya uydurma} olabilir. Bu nedenle "Güven ama Doğrula" (Trust but Verify) ilkesi gereği, çalışma gönderilmeden önce her referansın doğruluğu araştırmacı tarafından manuel doğrulanmalıdır.
    \item \textbf{Şeffaflık (Açıklama) ve Disiplin Kuralları:} Araştırma metodolojisinde yapay zekânın adı, sürüm bilgisi ve tam kullanım amacı şeffafça açıklanmalıdır. "Adalet ve Saygı" ilkesi (emeğe saygı) gereği yapay zekâ kullanımının gizlenmesi \textbf{akademik/disiplin suçu} olarak kabul edilir.
    \item \textbf{Entelektüel ve Yasal Yazar Sorumluluğu:} Elsevier, Springer ve saygın akademik dergilerin kurallarına göre, yapay zekâ araçları (ChatGPT vb.) bilimsel çalışmalarda \textbf{"ortak yazar" (co-author)} olarak listelenemez. Çünkü yapay zekânın tüzel veya ahlaki kişiliği yoktur, telif hakkı iddiasında bulunamaz, yasal sorumluluk ve hesap verebilirlik üstlenemez. Tüm yasal ve etik sorumluluk insan araştırmacıya aittir. Yapay zekanın entelektüel katkısının insan araştırmacı tarafından ikame edilmesi riskine \textbf{Özgünlük Erozyonu} denir.
    \item \textbf{Veri Anonimleştirme ve KVKK/GDPR Uyumluğu:} Şirket sırları, finansal veriler veya kişisel sağlık/kimlik bilgileri yapay zekâ platformlarına aktarılmadan önce kesinlikle temizlenmeli ve anonimleştirilmelidir. Aksi takdirde veri güvenliği ihlalleri oluşur.
    \item \textbf{Küresel Düzenlemeler ve EU AI Act:} Avrupa Birliği Yapay Zekâ Yasası, yapay zekâ sistemlerini taşıdıkları risk düzeylerine göre (Kabul Edilemez Risk, Yüksek Risk, Sınırlı Risk, Minimum Risk) sınıflandırarak denetleyen ve uygunsuzluk durumunda ağır yaptırımlar getiren ilk kapsamlı küresel regülasyondur.
\end{itemize}


\newpage
\chapter{Kelime Kartları ve Hızlı Tekrar}
Buradaki soru-cevap çiftleri, yapay zekâ ve istem mühendisliği konularındaki en kritik terimleri hızlıca pekiştirmeniz için tasarlanmıştır. Sınava girmeden önce bu kelime kartlarını gözden geçirmeniz tavsiye edilir.

\needspace{3.5cm}
\begin{tcolorbox}[colback=white, colframe=secondarycolor!50!black, arc=1.2mm, boxrule=0.8pt, top=8pt, bottom=8pt, left=12pt, right=12pt, segmentation style={dashed, secondarycolor!30, line width=0.8pt}]
    \noindent\textbf{\color{red!80!black}\sffamily Soru 1:} İstem (Prompt) mühendisliğinin temel amacı nedir?
\tcblower
    \noindent\textbf{\color{green!50!black}\sffamily Cevap:} Yapay zekâ ile etkileşimi tahminden çıkarıp, tekrarlanabilir bir mühendislik disiplinine dönüştürmektir.
\end{tcolorbox}
\vspace{6pt}
\needspace{3.5cm}
\begin{tcolorbox}[colback=white, colframe=secondarycolor!50!black, arc=1.2mm, boxrule=0.8pt, top=8pt, bottom=8pt, left=12pt, right=12pt, segmentation style={dashed, secondarycolor!30, line width=0.8pt}]
    \noindent\textbf{\color{red!80!black}\sffamily Soru 2:} İstem mühendisliğindeki 'Temel Denklem' nedir?
\tcblower
    \noindent\textbf{\color{green!50!black}\sffamily Cevap:} Girdi = Çıktı.
\end{tcolorbox}
\vspace{6pt}
\needspace{3.5cm}
\begin{tcolorbox}[colback=white, colframe=secondarycolor!50!black, arc=1.2mm, boxrule=0.8pt, top=8pt, bottom=8pt, left=12pt, right=12pt, segmentation style={dashed, secondarycolor!30, line width=0.8pt}]
    \noindent\textbf{\color{red!80!black}\sffamily Soru 3:} TÜBİTAK'a göre yapay zekâdan en verimli çıktıyı almak için ne yapılması gerekir?
\tcblower
    \noindent\textbf{\color{green!50!black}\sffamily Cevap:} Komut ve yönlendirmelerin bilinçli şekilde tasarlanması gerekir.
\end{tcolorbox}
\vspace{6pt}
\needspace{3.5cm}
\begin{tcolorbox}[colback=white, colframe=secondarycolor!50!black, arc=1.2mm, boxrule=0.8pt, top=8pt, bottom=8pt, left=12pt, right=12pt, segmentation style={dashed, secondarycolor!30, line width=0.8pt}]
    \noindent\textbf{\color{red!80!black}\sffamily Soru 4:} Bir istemin anatomisinde 'Sistem Talimatları' (System Instructions) neyi belirler?
\tcblower
    \noindent\textbf{\color{green!50!black}\sffamily Cevap:} Modelin benimseyeceği rolü, stili ve tonu belirler.
\end{tcolorbox}
\vspace{6pt}
\needspace{3.5cm}
\begin{tcolorbox}[colback=white, colframe=secondarycolor!50!black, arc=1.2mm, boxrule=0.8pt, top=8pt, bottom=8pt, left=12pt, right=12pt, segmentation style={dashed, secondarycolor!30, line width=0.8pt}]
    \noindent\textbf{\color{red!80!black}\sffamily Soru 5:} İstem bileşenlerinden hangisi modelin referans alacağı arka plan bilgisini sunar?
\tcblower
    \noindent\textbf{\color{green!50!black}\sffamily Cevap:} Bağlam (Contextual Info).
\end{tcolorbox}
\vspace{6pt}
\needspace{3.5cm}
\begin{tcolorbox}[colback=white, colframe=secondarycolor!50!black, arc=1.2mm, boxrule=0.8pt, top=8pt, bottom=8pt, left=12pt, right=12pt, segmentation style={dashed, secondarycolor!30, line width=0.8pt}]
    \noindent\textbf{\color{red!80!black}\sffamily Soru 6:} İstemde 'Görev' (Task) bileşeninin tanımı nedir?
\tcblower
    \noindent\textbf{\color{green!50!black}\sffamily Cevap:} Modelden beklenen temel işlem veya sorudur.
\end{tcolorbox}
\vspace{6pt}
\needspace{3.5cm}
\begin{tcolorbox}[colback=white, colframe=secondarycolor!50!black, arc=1.2mm, boxrule=0.8pt, top=8pt, bottom=8pt, left=12pt, right=12pt, segmentation style={dashed, secondarycolor!30, line width=0.8pt}]
    \noindent\textbf{\color{red!80!black}\sffamily Soru 7:} İstenen formatı gösteren doğru cevap şablonlarına ne ad verilir?
\tcblower
    \noindent\textbf{\color{green!50!black}\sffamily Cevap:} Örnekler (Few-shot Examples).
\end{tcolorbox}
\vspace{6pt}
\needspace{3.5cm}
\begin{tcolorbox}[colback=white, colframe=secondarycolor!50!black, arc=1.2mm, boxrule=0.8pt, top=8pt, bottom=8pt, left=12pt, right=12pt, segmentation style={dashed, secondarycolor!30, line width=0.8pt}]
    \noindent\textbf{\color{red!80!black}\sffamily Soru 8:} İteratif İyileştirme Döngüsü'nün ilk adımı nedir?
\tcblower
    \noindent\textbf{\color{green!50!black}\sffamily Cevap:} İhtiyacı Belirle (İstenen çıktıyı açıkça tanımla).
\end{tcolorbox}
\vspace{6pt}
\needspace{3.5cm}
\begin{tcolorbox}[colback=white, colframe=secondarycolor!50!black, arc=1.2mm, boxrule=0.8pt, top=8pt, bottom=8pt, left=12pt, right=12pt, segmentation style={dashed, secondarycolor!30, line width=0.8pt}]
    \noindent\textbf{\color{red!80!black}\sffamily Soru 9:} İteratif İyileştirme Döngüsü'nde komutun açık ve spesifik yazıldığı aşama hangisidir?
\tcblower
    \noindent\textbf{\color{green!50!black}\sffamily Cevap:} Yapılandır (3. Adım).
\end{tcolorbox}
\vspace{6pt}
\needspace{3.5cm}
\begin{tcolorbox}[colback=white, colframe=secondarycolor!50!black, arc=1.2mm, boxrule=0.8pt, top=8pt, bottom=8pt, left=12pt, right=12pt, segmentation style={dashed, secondarycolor!30, line width=0.8pt}]
    \noindent\textbf{\color{red!80!black}\sffamily Soru 10:} C.R.A.F.T. formülündeki 'C' harfi neyi temsil eder?
\tcblower
    \noindent\textbf{\color{green!50!black}\sffamily Cevap:} Context (Bağlam): Rol, hedef kitle ve ton.
\end{tcolorbox}
\vspace{6pt}
\needspace{3.5cm}
\begin{tcolorbox}[colback=white, colframe=secondarycolor!50!black, arc=1.2mm, boxrule=0.8pt, top=8pt, bottom=8pt, left=12pt, right=12pt, segmentation style={dashed, secondarycolor!30, line width=0.8pt}]
    \noindent\textbf{\color{red!80!black}\sffamily Soru 11:} C.R.A.F.T. formülündeki 'R' harfi (Request) neyi ifade eder?
\tcblower
    \noindent\textbf{\color{green!50!black}\sffamily Cevap:} Ulaşılması gereken spesifik hedefi ifade eder.
\end{tcolorbox}
\vspace{6pt}
\needspace{3.5cm}
\begin{tcolorbox}[colback=white, colframe=secondarycolor!50!black, arc=1.2mm, boxrule=0.8pt, top=8pt, bottom=8pt, left=12pt, right=12pt, segmentation style={dashed, secondarycolor!30, line width=0.8pt}]
    \textbf{Doğru Cevap: D}

    \vspace{2pt}
    Gem'lerin temel amacı, kullanıcının belirli görev ve roller için özelleştirilmiş talimatları bir kez tanımlayıp kaydetmesini sağlayarak, her sohbette aynı promptu yazma zorunluluğunu ortadan kaldırmaktır.
\end{tcolorbox}
\vspace{10pt}
% SORU 3
\needspace{5cm}
\vspace{8pt}
\noindent\textbf{\large\sffamily Soru 3.} \ OpenAI Custom GPTs oluştururken bilgi tabanı (knowledge base) için en fazla kaç dosya yüklenebilir?

\begin{enumerate}[label=\Alph*), leftmargin=25pt, topsep=4pt, itemsep=2pt]
        \item Sınırsız dosya
        \item 50 dosya
        \item 20 dosya
        \item 10 dosya
\end{enumerate}

\vspace{4pt}
\begin{tcolorbox}[title=Cevap ve Açıklama, colback=gray!3!white, colframe=primarycolor!90!black, arc=2mm, fonttitle=\bfseries\sffamily, top=6pt, bottom=6pt]
    \textbf{Doğru Cevap: C}

    \vspace{2pt}
    OpenAI platformunda bir Custom GPT oluştururken bilgi tabanına (Knowledge Base) yüklenebilecek maksimum dosya sayısı 20 olarak sınırlandırılmıştır.
\end{tcolorbox}
\vspace{10pt}
% SORU 4
\needspace{5cm}
\vspace{8pt}
\noindent\textbf{\large\sffamily Soru 4.} \ Aşağıdaki platformlardan hangisi ücretsiz olarak özelleştirilmiş AI asistanı oluşturmaya olanak tanır?

\begin{enumerate}[label=\Alph*), leftmargin=25pt, topsep=4pt, itemsep=2pt]
        \item Claude Projects
        \item OpenAI Custom GPTs
        \item GPT Store
        \item Google Gemini Gems
\end{enumerate}

\vspace{4pt}
\begin{tcolorbox}[title=Cevap ve Açıklama, colback=gray!3!white, colframe=primarycolor!90!black, arc=2mm, fonttitle=\bfseries\sffamily, top=6pt, bottom=6pt]
    \textbf{Doğru Cevap: B}

    \vspace{2pt}
    Gelişmiş asistan oluşturma özellikleri (Claude Projects, OpenAI Custom GPTs ve Google Gemini Gems) normal şartlarda ilgili platformların ücretli aboneliklerini (Claude Pro, ChatGPT Plus, Gemini Advanced) gerektirir. Ancak OpenAI, Mayıs 2024 güncellemesiyle daha önceden oluşturulmuş Custom GPT'lerin kullanımını ücretsiz kullanıcılara açmıştır. Bu doğrultuda soruda hedeflenen en yakın seçenek B'dir.
\end{tcolorbox}
\vspace{10pt}
% SORU 5
\needspace{5cm}
\vspace{8pt}
\noindent\textbf{\large\sffamily Soru 5.} \ Bağlam (context) penceresi bakımından en geniş kapasiteyi hangi platform sunmaktadır?

\begin{enumerate}[label=\Alph*), leftmargin=25pt, topsep=4pt, itemsep=2pt]
        \item Google Gemini Gems (1M token)
        \item Claude Projects (200K token)
        \item Tüm platformlar 128K token sunar
        \item OpenAI Custom GPTs (128K token)
\end{enumerate}

\vspace{4pt}
\begin{tcolorbox}[title=Cevap ve Açıklama, colback=gray!3!white, colframe=primarycolor!90!black, arc=2mm, fonttitle=\bfseries\sffamily, top=6pt, bottom=6pt]
    \textbf{Doğru Cevap: A}

    \vspace{2pt}
    Google Gemini, 1 milyon token (ve hatta bazı modellerde 2 milyona kadar) destekleyen devasa bağlam (context) penceresi kapasitesiyle rakiplerine göre en geniş veri işleme hacmini sunmaktadır.
\end{tcolorbox}
\vspace{10pt}
% SORU 6
\needspace{5cm}
\vspace{8pt}
\noindent\textbf{\large\sffamily Soru 6.} \ İyi bir asistan oluşturmak için hazırlanan prompt şablonunda 'KURALLAR' bölümünün amacı nedir?

\begin{enumerate}[label=\Alph*), leftmargin=25pt, topsep=4pt, itemsep=2pt]
        \item Asistanın neyi yapması ve kesinlikle neyi yapmaması gerektiğini belirlemek
        \item Asistanın sadece hangi dilde konuşacağını seçmek
        \item Yüklenecek dosyaların formatını denetlemek
        \item Sadece asistanın ismini belirlemek
\end{enumerate}

\vspace{4pt}
\begin{tcolorbox}[title=Cevap ve Açıklama, colback=gray!3!white, colframe=primarycolor!90!black, arc=2mm, fonttitle=\bfseries\sffamily, top=6pt, bottom=6pt]
    \textbf{Doğru Cevap: A}

    \vspace{2pt}
    Prompt şablonundaki 'KURALLAR' veya 'TALİMATLAR' bölümü, asistanın davranış sınırlarını çizer. Asistanın hangi görevleri yerine getireceğini, hangi üslubu kullanacağını ve özellikle neleri yapmaktan kaçınması gerektiğini (kısıtlamaları) netleştirir.
\end{tcolorbox}
\vspace{10pt}
% SORU 7
\needspace{5cm}
\vspace{8pt}
\noindent\textbf{\large\sffamily Soru 7.} \ Bir startup ekibi için 'Custom GPTs' platformunun 'Gems'e göre en büyük avantajı nedir?

\begin{enumerate}[label=\Alph*), leftmargin=25pt, topsep=4pt, itemsep=2pt]
        \item Sadece Google Drive dosyalarıyla çalışması
        \item Ücretsiz olması sayesinde maliyet düşürmesi
        \item Ekiple paylaşılabilir olması ve API entegrasyonu sunması
        \item Prompt yazmaya ihtiyaç duymaması
\end{enumerate}

\vspace{4pt}
\begin{tcolorbox}[title=Cevap ve Açıklama, colback=gray!3!white, colframe=primarycolor!90!black, arc=2mm, fonttitle=\bfseries\sffamily, top=6pt, bottom=6pt]
    \textbf{Doğru Cevap: C}

    \vspace{2pt}
    OpenAI Custom GPTs, API entegrasyonları (Actions) ve ekiple/kurum içi paylaşım imkanı sunarak startup ekiplerinin kendi sistemleriyle entegre çalışabilen özelleştirilmiş botlar geliştirmesine daha elverişlidir.
\end{tcolorbox}
\vspace{10pt}
% SORU 8
\needspace{5cm}
\vspace{8pt}
\noindent\textbf{\large\sffamily Soru 8.} \ Özelleştirilmiş bir AI oluştururken yapılan en yaygın hata aşağıdakilerden hangisidir?

\begin{enumerate}[label=\Alph*), leftmargin=25pt, topsep=4pt, itemsep=2pt]
        \item 'Yardımcı ol' gibi çok genel ve belirsiz talimatlar kullanmak
        \item Bilgi tabanına örnek dosyalar eklemek
        \item Asistanı yayınlamadan önce test etmek
        \item Çok spesifik talimatlar vermek
\end{enumerate}

\vspace{4pt}
\begin{tcolorbox}[title=Cevap ve Açıklama, colback=gray!3!white, colframe=primarycolor!90!black, arc=2mm, fonttitle=\bfseries\sffamily, top=6pt, bottom=6pt]
    \textbf{Doğru Cevap: A}

    \vspace{2pt}
    Yapay zeka asistanı oluştururken en yaygın hata, "Kullanıcıya her konuda yardımcı ol" veya "İyi bir asistan ol" gibi yoruma açık, belirsiz ve çok genel talimatlar vermektir. Bu durum asistanın tutarsız ve verimsiz çalışmasına yol açar.
\end{tcolorbox}
\vspace{10pt}
% SORU 9
\needspace{5cm}
\vspace{8pt}
\noindent\textbf{\large\sffamily Soru 9.} \ Claude Projects hakkında dokümanda verilen bilgilerden yola çıkarak, paylaşım özelliği için ne söylenebilir?

\begin{enumerate}[label=\Alph*), leftmargin=25pt, topsep=4pt, itemsep=2pt]
        \item Google Gems gibi hiç paylaşım yapılamaz
        \item Sadece link ile dünya çapında paylaşılabilir
        \item Paylaşım seçenekleri sınırlıdır
        \item Tamamen herkese açık bir marketplace'e sahiptir
\end{enumerate}

\vspace{4pt}
\begin{tcolorbox}[title=Cevap ve Açıklama, colback=gray!3!white, colframe=primarycolor!90!black, arc=2mm, fonttitle=\bfseries\sffamily, top=6pt, bottom=6pt]
    \textbf{Doğru Cevap: C}

    \vspace{2pt}
    Claude Projects'te projeler yalnızca organizasyon (Team) içindeki diğer kullanıcılarla paylaşılabilir. ChatGPT'deki gibi genel kullanıma açık bir linkle dış dünyayla paylaşma veya GPT Store benzeri herkese açık bir pazaryeri (marketplace) bulunmamaktadır. Bu yüzden paylaşım seçenekleri sınırlıdır.
\end{tcolorbox}
\vspace{10pt}
% SORU 10
\needspace{5cm}
\vspace{8pt}
\noindent\textbf{\large\sffamily Soru 10.} \ Aşağıdaki senaryolardan hangisi için Google Gemini Gems kullanmak en mantıklıdır?

\begin{enumerate}[label=\Alph*), leftmargin=25pt, topsep=4pt, itemsep=2pt]
        \item Harici bir stok takip yazılımına API ile bağlanan bir bot
        \item Asistanı GPT Store'da satarak gelir elde etmek
        \item Müşteri destek botunu tüm dünya ile paylaşmak
        \item Drive'daki ders notlarına erişip özet çıkaran bir öğrenci asistanı
\end{enumerate}

\vspace{4pt}
\begin{tcolorbox}[title=Cevap ve Açıklama, colback=gray!3!white, colframe=primarycolor!90!black, arc=2mm, fonttitle=\bfseries\sffamily, top=6pt, bottom=6pt]
    \textbf{Doğru Cevap: D}

    \vspace{2pt}
    Google Gemini Gems, Google Workspace (Google Drive, Gmail, Docs) araçlarıyla yerleşik entegrasyona sahiptir. Dolayısıyla Google Drive'daki notlara doğrudan erişip bunlardan özet çıkarma senaryosu Gemini Gems için en uygun ve verimli kullanım alanıdır.
\end{tcolorbox}
\vspace{10pt}

\section{Gem Testi}
% SORU 11
\needspace{5cm}
\vspace{8pt}
\noindent\textbf{\large\sffamily Soru 11.} \ Gem mimarisinin temel amacı aşağıdakilerden hangisidir?

\begin{enumerate}[label=\Alph*), leftmargin=25pt, topsep=4pt, itemsep=2pt]
        \item Sadece metin tabanlı yanıtların hızını artırmak.
        \item Yapay zekanın genel bilgi birikimini tamamen sıfırlamak.
        \item İnternet bağlantısı olmadan Gemini kullanımını sağlamak.
        \item Her sohbette yeni bir istem (prompt) yazma zorunluluğunu ortadan kaldırmak.
\end{enumerate}

\vspace{4pt}
\begin{tcolorbox}[title=Cevap ve Açıklama, colback=gray!3!white, colframe=primarycolor!90!black, arc=2mm, fonttitle=\bfseries\sffamily, top=6pt, bottom=6pt]
    \textbf{Doğru Cevap: D}
    
    \vspace{2pt}
    Gem'ler (Gems), kullanıcının belirli görevler veya roller için özelleştirilmiş talimatları (sistem istemlerini) bir kez tanımlayıp kaydetmesini sağlar. Böylece her yeni sohbette aynı talimatları tekrar yazma zorunluluğu ortadan kalkar.
\end{tcolorbox}
\vspace{10pt}
% SORU 12
\needspace{5cm}
\vspace{8pt}
\noindent\textbf{\large\sffamily Soru 12.} \ Gem kurulum sürecinde 'Laboratuvar' olarak adlandırılan aşamalardan hangisi sistemin aktifleşmesini sağlar?

\begin{enumerate}[label=\Alph*), leftmargin=25pt, topsep=4pt, itemsep=2pt]
        \item Keşfet
        \item Test Et
        \item Kimlik Ver
        \item Sisteme Ekle
\end{enumerate}

\vspace{4pt}
\begin{tcolorbox}[title=Cevap ve Açıklama, colback=gray!3!white, colframe=primarycolor!90!black, arc=2mm, fonttitle=\bfseries\sffamily, top=6pt, bottom=6pt]
    \textbf{Doğru Cevap: D}

    \vspace{2pt}
    Gem oluşturma sürecinde yapılandırma tamamlandıktan sonra asistanı kaydetmek ve kullanıma sunmak (aktifleşmesini sağlamak) için 'Sisteme Ekle' seçeneği kullanılır.
\end{tcolorbox}
\vspace{10pt}
% SORU 13
\needspace{5cm}
\vspace{8pt}
\noindent\textbf{\large\sffamily Soru 13.} \ Gem'in hangi rolü oynayacağını ve nasıl bir persona çizeceğini belirleyen yapı taşı hangisidir?

\begin{enumerate}[label=\Alph*), leftmargin=25pt, topsep=4pt, itemsep=2pt]
        \item Görev
        \item Karakter
        \item Bağlam
        \item Biçim
\end{enumerate}

\vspace{4pt}
\begin{tcolorbox}[title=Cevap ve Açıklama, colback=gray!3!white, colframe=primarycolor!90!black, arc=2mm, fonttitle=\bfseries\sffamily, top=6pt, bottom=6pt]
    \textbf{Doğru Cevap: B}

    \vspace{2pt}
    Gem'in veya promptun hangi kimliğe bürüneceğini, nasıl bir ses tonu ve üslup kullanacağını (personayı) belirleyen yapı taşı 'Karakter'dir.
\end{tcolorbox}
\vspace{10pt}
% SORU 14
\needspace{5cm}
\vspace{8pt}
\noindent\textbf{\large\sffamily Soru 14.} \ Bir Gem tasarlarken çıktının madde işaretleri veya tablolar şeklinde olmasını istiyorsanız hangi sütuna ayrıntı eklemelisiniz?

\begin{enumerate}[label=\Alph*), leftmargin=25pt, topsep=4pt, itemsep=2pt]
        \item Karakter
        \item Görev
        \item Biçim
        \item Bağlam
\end{enumerate}

\vspace{4pt}
\begin{tcolorbox}[title=Cevap ve Açıklama, colback=gray!3!white, colframe=primarycolor!90!black, arc=2mm, fonttitle=\bfseries\sffamily, top=6pt, bottom=6pt]
    \textbf{Doğru Cevap: C}

    \vspace{2pt}
    Çıktının görsel yapısını, listeler, tablolar, kod blokları veya paragraflar halinde mi olacağını belirleyen prompt bileşeni 'Biçim' (Format) seçeneğidir.
\end{tcolorbox}
\vspace{10pt}
% SORU 15
\needspace{5cm}
\vspace{8pt}
\noindent\textbf{\large\sffamily Soru 15.} \ Beyin Fırtınası Yardımcısı (Fikir Üretim Motoru) olarak tasarlanan bir Gem'de hangi sütunların vurgusu en yüksektir?

\begin{enumerate}[label=\Alph*), leftmargin=25pt, topsep=4pt, itemsep=2pt]
        \item Karakter ve Biçim
        \item Görev ve Biçim
        \item Sadece Görev
        \item Karakter ve Bağlam
\end{enumerate}

\vspace{4pt}
\begin{tcolorbox}[title=Cevap ve Açıklama, colback=gray!3!white, colframe=primarycolor!90!black, arc=2mm, fonttitle=\bfseries\sffamily, top=6pt, bottom=6pt]
    \textbf{Doğru Cevap: B}

    \vspace{2pt}
    Beyin fırtınası yardımcısında fikirlerin listelenmesi ve yapılandırılması için en kritik bileşenler yerine getirilecek 'Görev' ve çıktının yapısını belirleyen 'Biçim'dir.
\end{tcolorbox}
\vspace{10pt}
% SORU 16
\needspace{5cm}
\vspace{8pt}
\noindent\textbf{\large\sffamily Soru 16.} \ Gem'inize 'Hafıza' eklemek için kullanılan yöntem aşağıdakilerden hangisidir?

\begin{enumerate}[label=\Alph*), leftmargin=25pt, topsep=4pt, itemsep=2pt]
        \item Gemini'ye önceki sohbetleri hatırlatmasını söylemek.
        \item Karakter sütununa daha fazla cümle yazmak.
        \item Talimat Döngüsünü daha fazla tekrarlamak.
        \item Harici dosya (PDF, Drive vb.) yükleme.
\end{enumerate}

\vspace{4pt}
\begin{tcolorbox}[title=Cevap ve Açıklama, colback=gray!3!white, colframe=primarycolor!90!black, arc=2mm, fonttitle=\bfseries\sffamily, top=6pt, bottom=6pt]
    \textbf{Doğru Cevap: D}

    \vspace{2pt}
    Gem'lere bilgi tabanı (hafıza) eklemek ve kendi verilerimiz üzerinde çalışmasını sağlamak için harici dosya (PDF, TXT, Google Drive vb.) yükleme yöntemi kullanılır.
\end{tcolorbox}
\vspace{10pt}
% SORU 17
\needspace{5cm}
\vspace{8pt}
\noindent\textbf{\large\sffamily Soru 17.} \ Gem'in yanıt verirken bilgi alıntılarından gizlemesini sağlayan özellik hangisidir?

\begin{enumerate}[label=\Alph*), leftmargin=25pt, topsep=4pt, itemsep=2pt]
        \item Hafıza Temizleyici
        \item Görünmez Kaynaklar
        \item Alıntı Filtresi
        \item Bağlam Kilidi
\end{enumerate}

\vspace{4pt}
\begin{tcolorbox}[title=Cevap ve Açıklama, colback=gray!3!white, colframe=primarycolor!90!black, arc=2mm, fonttitle=\bfseries\sffamily, top=6pt, bottom=6pt]
    \textbf{Doğru Cevap: C}

    \vspace{2pt}
    Yapay zekanın yüklenen dosyalardan alıntı göstermesini engelleyen veya gizleyen ayara 'Alıntı Filtresi' denir.
\end{tcolorbox}
\vspace{10pt}
% SORU 18
\needspace{5cm}
\vspace{8pt}
\noindent\textbf{\large\sffamily Soru 18.} \ Mükemmel Gem'in Formülü (Sentez) aşağıdakilerden hangisinin birleşimidir?

\begin{enumerate}[label=\Alph*), leftmargin=25pt, topsep=4pt, itemsep=2pt]
        \item 4 Temel Sütun + Özel Bilgi Havuzu + Talimat Döngüsü
        \item Karakter + Görev + Bağlam + Biçim
        \item Test Et + Kaydet + Paylaş
        \item İstem (Prompt) + Analiz + Çıktı
\end{enumerate}

\vspace{4pt}
\begin{tcolorbox}[title=Cevap ve Açıklama, colback=gray!3!white, colframe=primarycolor!90!black, arc=2mm, fonttitle=\bfseries\sffamily, top=6pt, bottom=6pt]
    \textbf{Doğru Cevap: A}
    
    \vspace{2pt}
    Sentez formülü; prompt tasarımının yapı taşları olan \textbf{4 Temel Sütunu} (Karakter, Görev, Bağlam, Biçim), Gem'in kullanacağı \textbf{Özel Bilgi Havuzunu} (Knowledge Base) ve sürekli geri bildirimle iyileştirme sağlayan \textbf{Talimat Döngüsünü} kapsar.
\end{tcolorbox}
\vspace{10pt}
% SORU 19
\needspace{5cm}
\vspace{8pt}
\noindent\textbf{\large\sffamily Soru 19.} \ Kodlama Yardımcısı olarak tasarlanan bir Gem'in 'Biçim' (Format) beklentisi aşağıdakilerden hangisidir?

\begin{enumerate}[label=\Alph*), leftmargin=25pt, topsep=4pt, itemsep=2pt]
        \item Kopyalanabilir bloklar halinde ve değişken açıklamalarıyla sunum.
        \item Şiirsel ve coşkulu bir üslup.
        \item En az 3 farklı programlama dilinde çıktı.
        \item Sadece metin üzerinden teorik anlatım.
\end{enumerate}

\vspace{4pt}
\begin{tcolorbox}[title=Cevap ve Açıklama, colback=gray!3!white, colframe=primarycolor!90!black, arc=2mm, fonttitle=\bfseries\sffamily, top=6pt, bottom=6pt]
    \textbf{Doğru Cevap: A}
    
    \vspace{2pt}
    Bir kodlama yardımcısının en verimli çıktısı, kullanıcının doğrudan kopyalayıp editörüne yapıştırabileceği kod blokları (markdown syntax) sunması ve bu koddaki değişkenleri/fonksiyonları anlaşılır açıklamalarla detaylandırmasıdır.
\end{tcolorbox}
\vspace{10pt}

\newpage
\section{İstem Testi}
% SORU 20
\needspace{5cm}
\vspace{8pt}
\noindent\textbf{\large\sffamily Soru 20.} \ İstem (Prompt) mühendisliğinin temel amacı aşağıdakilerden hangisidir?

\begin{enumerate}[label=\Alph*), leftmargin=25pt, topsep=4pt, itemsep=2pt]
        \item Yapay zekâ ile etkileşimi tahminden çıkarıp tekrarlanabilir bir disipline dönüştürmek.
        \item Modelin eğitim verisindeki tüm bilgileri internete sızdırmak.
        \item Yapay zekâ modellerini tamamen ortadan kaldırıp geleneksel kodlamaya dönmek.
        \item Yapay zekânın işlemci hızını donanımsal olarak artırmak.
\end{enumerate}

\vspace{4pt}
\begin{tcolorbox}[title=Cevap ve Açıklama, colback=gray!3!white, colframe=primarycolor!90!black, arc=2mm, fonttitle=\bfseries\sffamily, top=6pt, bottom=6pt]
    \textbf{Doğru Cevap: A}

    \vspace{2pt}
    İstem Mühendisliğinin temel amacı, yapay zekayla etkileşimi rastgele denemelerden çıkarıp sistemli, öngörülebilir ve tekrarlanabilir bir disipline dönüştürmektir.
\end{tcolorbox}
\vspace{10pt}
% SORU 21
\needspace{5cm}
\vspace{8pt}
\noindent\textbf{\large\sffamily Soru 21.} \ İstem Mühendisliğindeki 'Temel Denklem' aşağıdakilerden hangisidir?

\begin{enumerate}[label=\Alph*), leftmargin=25pt, topsep=4pt, itemsep=2pt]
        \item Donanım + Yazılım = YapayZekâ
        \item Hız x Veri = Doğruluk
        \item Girdi = Çıktı
        \item Prompt - Bağlam = Yanıt
\end{enumerate}

\vspace{4pt}
\begin{tcolorbox}[title=Cevap ve Açıklama, colback=gray!3!white, colframe=primarycolor!90!black, arc=2mm, fonttitle=\bfseries\sffamily, top=6pt, bottom=6pt]
    \textbf{Doğru Cevap: C}

    \vspace{2pt}
    İstem mühendisliğindeki temel denklem 'Girdi = Çıktı' (GIGO - Garbage In, Garbage Out) prensibidir; modele ne kadar kaliteli girdi sağlanırsa o kadar kaliteli çıktı alınır.
\end{tcolorbox}
\vspace{10pt}
% SORU 22
\needspace{5cm}
\vspace{8pt}
\noindent\textbf{\large\sffamily Soru 22.} \ C.R.A.F.T. oluşturma formülündeki 'A' (Actions) harfi neyi temsil eder?

\begin{enumerate}[label=\Alph*), leftmargin=25pt, topsep=4pt, itemsep=2pt]
        \item Kullanılacak olan yapay zekâ modelinin adı.
        \item Modelin izlemesi gereken adım adım talimatlar.
        \item Yapay zekânın hata yapma olasılığı.
        \item Analitik verilerin tabloya dönüştürmesi.
\end{enumerate}

\vspace{4pt}
\begin{tcolorbox}[title=Cevap ve Açıklama, colback=gray!3!white, colframe=primarycolor!90!black, arc=2mm, fonttitle=\bfseries\sffamily, top=6pt, bottom=6pt]
    \textbf{Doğru Cevap: B}

    \vspace{2pt}
    C.R.A.F.T. formülünde 'A' (Actions - Eylemler) harfi, yapay zekanın görevi yerine getirirken izlemesi gereken adım adım talimatları ve gerçekleştireceği eylemleri temsil eder.
\end{tcolorbox}
\vspace{10pt}
% SORU 23
\needspace{5cm}
\vspace{8pt}
\noindent\textbf{\large\sffamily Soru 23.} \ C.L.E.A.R. kalite kontrol listesine göre 'Explicit' (Açık) kriteri ne anlama gelir?

\begin{enumerate}[label=\Alph*), leftmargin=25pt, topsep=4pt, itemsep=2pt]
        \item Modelin internete bağlı olup olmadığının kontrolü.
        \item İstemin mümkün olduğunca karmaşık ve uzun yazılması.
        \item Cevabın sadece İngilizce dilinde verilmesi.
        \item Çıktı yönlendirmelerinin 'Kısa cevap ver' yerine '50 kelimeyi geçme' gibi net olması.
\end{enumerate}

\vspace{4pt}
\begin{tcolorbox}[title=Cevap ve Açıklama, colback=gray!3!white, colframe=primarycolor!90!black, arc=2mm, fonttitle=\bfseries\sffamily, top=6pt, bottom=6pt]
    \textbf{Doğru Cevap: D}

    \vspace{2pt}
    C.L.E.A.R. kalite kontrol listesindeki 'Explicit' (Açık/Net) kriteri, yapay zekaya verilen çıktı komutlarının ve sınırlamaların (örneğin '50 kelimeyi geçme') net ve ölçülebilir olması gerektiğini belirtir.
\end{tcolorbox}
\vspace{10pt}
% SORU 24
\needspace{5cm}
\vspace{8pt}
\noindent\textbf{\large\sffamily Soru 24.} \ Modele bir Persona (Rol) atamanın en önemli avantajlarından biri nedir?

\begin{enumerate}[label=\Alph*), leftmargin=25pt, topsep=4pt, itemsep=2pt]
        \item Halüsinasyon (uydurma) oranını düşürmesi ve göreve özgü terminolojiyi aktive etmesi.
        \item Tüm güvenlik filtrelerini devre dışı bırakması.
        \item Yapay zekânın insani duygular kazanmasını sağlaması.
        \item Modelin daha hızlı yanıt vermesini sağlaması.
\end{enumerate}

\vspace{4pt}
\begin{tcolorbox}[title=Cevap ve Açıklama, colback=gray!3!white, colframe=primarycolor!90!black, arc=2mm, fonttitle=\bfseries\sffamily, top=6pt, bottom=6pt]
    \textbf{Doğru Cevap: A}

    \vspace{2pt}
    Modele belirli bir rol veya persona atamak, yapay zekanın o konuya özgü terminolojiyi ve üslubu kullanmasını tetiklerken, aynı zamanda alakasız bilgiler üretmesini (halüsinasyonu) azaltır.
\end{tcolorbox}
\vspace{10pt}
% SORU 25
\needspace{5cm}
\vspace{8pt}
\noindent\textbf{\large\sffamily Soru 25.} \ İstem Olgunluk Eğrisi'ne (Maturity Curve) göre, modelden karmaşık mantık yürütmesini beklediğimiz en üst seviye teknik hangisidir?

\begin{enumerate}[label=\Alph*), leftmargin=25pt, topsep=4pt, itemsep=2pt]
        \item Düşünce Zinciri (Chain of Thought - CoT)
        \item Az Örnekli İstem (Few-Shot)
        \item Doğrudan Komut (Zero-Shot)
        \item Karakter Sınırlama
\end{enumerate}

\vspace{4pt}
\begin{tcolorbox}[title=Cevap ve Açıklama, colback=gray!3!white, colframe=primarycolor!90!black, arc=2mm, fonttitle=\bfseries\sffamily, top=6pt, bottom=6pt]
    \textbf{Doğru Cevap: A}

    \vspace{2pt}
    İstem Olgunluk Eğrisi'nin en üst düzey tekniklerinden biri olan Düşünce Zinciri (Chain of Thought), modelin karmaşık problemleri çözerken adım adım akıl yürüterek doğru sonuca ulaşmasını kolaylaştırır.
\end{tcolorbox}
\vspace{10pt}
% SORU 26
\needspace{5cm}
\vspace{8pt}
\noindent\textbf{\large\sffamily Soru 26.} \ Few-Shot (Az Örnekli) istemleme tekniğinin temel farkı nedir?

\begin{enumerate}[label=\Alph*), leftmargin=25pt, topsep=4pt, itemsep=2pt]
        \item AI'ya sadece tek bir kelime söyleyip geri kalanını tahmin etmesini beklemek.
        \item Modelin eğitim verisindeki tüm örnekleri silmek.
        \item AI'ya örnek vermeden sadece 'Lütfen' diyerek ricada bulunmak.
        \item AI'ya istenen formatı tarif etmek yerine birkaç örnekle göstermek.
\end{enumerate}

\vspace{4pt}
\begin{tcolorbox}[title=Cevap ve Açıklama, colback=gray!3!white, colframe=primarycolor!90!black, arc=2mm, fonttitle=\bfseries\sffamily, top=6pt, bottom=6pt]
    \textbf{Doğru Cevap: D}
    
    \vspace{2pt}
    Few-shot (Az Örnekli) istemleme, modele ne yapacağını sadece tarif etmek yerine girdilerin ve istenen çıktıların birkaç örneğini göstererek format ve ton uyumunu artırmak için kullanılır.
\end{tcolorbox}
\vspace{10pt}
% SORU 27
\needspace{5cm}
\vspace{8pt}
\noindent\textbf{\large\sffamily Soru 27.} \ İleri düzey bir taktik olan 'Chunking' (Parçalara Bölme) neden gereklidir?

\begin{enumerate}[label=\Alph*), leftmargin=25pt, topsep=4pt, itemsep=2pt]
        \item Yapay zekânın daha fazla elektrik tüketmesini sağlamak için.
        \item Metni daha zor okunur hale getirmek için.
        \item Büyük görevlerin tek seferde verilmesi durumunda sistemin boğulmasını ve sığ cevaplar üretmesini engellemek için.
        \item Modelin belleğini tamamen sıfırlamak için.
\end{enumerate}

\vspace{4pt}
\begin{tcolorbox}[title=Cevap ve Açıklama, colback=gray!3!white, colframe=primarycolor!90!black, arc=2mm, fonttitle=\bfseries\sffamily, top=6pt, bottom=6pt]
    \textbf{Doğru Cevap: C}

    \vspace{2pt}
    Chunking (Parçalara Bölme), çok büyük veya karmaşık görevlerin tek seferde verilip modelin bağlam sınırlarını zorlayarak sığ veya eksik cevaplar üretmesini engellemek için görevi küçük parçalara bölme işlemidir.
\end{tcolorbox}
\vspace{10pt}
% SORU 28
\needspace{5cm}
\vspace{8pt}
\noindent\textbf{\large\sffamily Soru 28.} \ Self-Consistency (Öz-Tutarlılık) tekniği hangi durumlarda en etkilidir?

\begin{enumerate}[label=\Alph*), leftmargin=25pt, topsep=4pt, itemsep=2pt]
        \item Modelin sürekli aynı şakayı yapmasını sağlamak için.
        \item Kritik kararlar, matematik problemleri veya belirsizlik içeren durumlarda en tutarlı cevabı seçmek için.
        \item AI'nın internetten rastgele resimler indirmesi için.
        \item Kullanıcının şifrelerini hatırlatması için.
\end{enumerate}

\vspace{4pt}
\begin{tcolorbox}[title=Cevap ve Açıklama, colback=gray!3!white, colframe=primarycolor!90!black, arc=2mm, fonttitle=\bfseries\sffamily, top=6pt, bottom=6pt]
    \textbf{Doğru Cevap: B}

    \vspace{2pt}
    Self-Consistency (Öz-Tutarlılık), modelin aynı soruya birden fazla farklı düşünce yoluyla yanıtlar üretip aralarından en tutarlı ve en çok tekrarlanan yanıtı seçerek doğruluk oranını artırma tekniğidir.
\end{tcolorbox}
\vspace{10pt}
% SORU 29
\needspace{5cm}
\vspace{8pt}
\noindent\textbf{\large\sffamily Soru 29.} \ Sentez Matrisi'ne göre 'Denetçinin Kontrol Listesi' hangisidir ve amacı nedir?

\begin{enumerate}[label=\Alph*), leftmargin=25pt, topsep=4pt, itemsep=2pt]
        \item CRAFTED listesidir; amacı ilk talimatı kusursuzca inşa etmektir.
        \item CLEAR listesidir; amacı çıktıyı denetlemek ve komutu rafine etmektir.
        \item JSON formatıdır; amacı veriyi tabloya dönüştürmektir.
        \item Python kodudur; amacı algoritmayı çalıştırmaktır.
\end{enumerate}

\vspace{4pt}
\begin{tcolorbox}[title=Cevap ve Açıklama, colback=gray!3!white, colframe=primarycolor!90!black, arc=2mm, fonttitle=\bfseries\sffamily, top=6pt, bottom=6pt]
    \textbf{Doğru Cevap: B}

    \vspace{2pt}
    Sentez Matrisinde 'Denetçinin Kontrol Listesi' CLEAR kontrol listesidir. Bu liste, üretilen çıktıyı ve komutları denetleyip en verimli şekilde rafine etmek için kullanılır.
\end{tcolorbox}
\vspace{10pt}
% SORU 30
\needspace{5cm}
\vspace{8pt}
\noindent\textbf{\large\sffamily Soru 30.} \ İstem (Prompt) Mühendisliği mimarisindeki 'Temel Denklem' aşağıdakilerden hangisidir?

\begin{enumerate}[label=\Alph*), leftmargin=25pt, topsep=4pt, itemsep=2pt]
        \item Yapay Zekâ + Veri = Bilgi
        \item Girdi = Çıktı
        \item Komut x Bağlam = Tahmin
        \item Algoritma = Sonuç
\end{enumerate}

\vspace{4pt}
\begin{tcolorbox}[title=Cevap ve Açıklama, colback=gray!3!white, colframe=primarycolor!90!black, arc=2mm, fonttitle=\bfseries\sffamily, top=6pt, bottom=6pt]
    \textbf{Doğru Cevap: B}

    \vspace{2pt}
    İstem Mühendisliği mimarisindeki temel denklem, 'Girdi = Çıktı' (Input = Output) şeklinde ifade edilerek girdinin çıktıyı doğrudan belirlediğini vurgular.
\end{tcolorbox}
\vspace{10pt}
% SORU 31
\needspace{5cm}
\vspace{8pt}
\noindent\textbf{\large\sffamily Soru 31.} \ CRAFTED formülündeki 'F' harfi neyi temsil eder ve ne amaçla kullanılır?

\begin{enumerate}[label=\Alph*), leftmargin=25pt, topsep=4pt, itemsep=2pt]
        \item Format (Biçim): Çıktının hangi yapıda olacağını belirlemek için kullanılır.
        \item Function (İşlev): Modelin gerçekleştireceği ana eylemi tanımlar.
        \item Frame (Sınurlar): Nelerin dışarıda bırakılacağını belirlemek için kullanılır.
        \item Feedback (Geri Bildirim): Modelin yanıtını iyileştirmek için kullanılır.
\end{enumerate}

\vspace{4pt}
\begin{tcolorbox}[title=Cevap ve Açıklama, colback=gray!3!white, colframe=primarycolor!90!black, arc=2mm, fonttitle=\bfseries\sffamily, top=6pt, bottom=6pt]
    \textbf{Doğru Cevap: A}

    \vspace{2pt}
    CRAFTED formülünde 'F' harfi 'Format' (Biçim) anlamına gelir ve yapay zekanın üreteceği çıktının görsel yapısını, listeler, tablolar veya kod blokları halinde olacağını belirler.
\end{tcolorbox}
\vspace{10pt}
% SORU 32
\needspace{5cm}
\vspace{8pt}
\noindent\textbf{\large\sffamily Soru 32.} \ Modele bir 'Rol' veya 'Persona' atamanın temel faydalarından biri aşağıdakilerden hangisidir?

\begin{enumerate}[label=\Alph*), leftmargin=25pt, topsep=4pt, itemsep=2pt]
        \item Modelin halüsinasyon (uydurma) oranını düşürmesi.
        \item Modelin internetten güncel veri çekmesini sağlaması.
        \item Modelin daha az kelime kullanarak yanıt vermesini sağlaması.
        \item Modelin işlem hızını (token/saniye) artırması.
\end{enumerate}

\vspace{4pt}
\begin{tcolorbox}[title=Cevap ve Açıklama, colback=gray!3!white, colframe=primarycolor!90!black, arc=2mm, fonttitle=\bfseries\sffamily, top=6pt, bottom=6pt]
    \textbf{Doğru Cevap: A}

    \vspace{2pt}
    Yapay zeka modellerine bir 'Rol' veya 'Persona' atandığında, model o göreve odaklanır ve rastgele veya alakasız bilgiler uydurma (halüsinasyon) oranı önemli ölçüde düşer.
\end{tcolorbox}
\vspace{10pt}
% SORU 33
\needspace{5cm}
\vspace{8pt}
\noindent\textbf{\large\sffamily Soru 33.} \ Few-Shot Prompting (Az Örnekli İstem) tekniği en çok hangi durumda tercih edilir?

\begin{enumerate}[label=\Alph*), leftmargin=25pt, topsep=4pt, itemsep=2pt]
        \item Modelin belleğinin tamamen temizlenmesi gerektiğinde.
        \item Çok basit ve doğrudan bir çeviri görevi yapıldığında.
        \item Karmaşık bir matematik problemini tek adımda çözmek için.
        \item Modelin belirli bir formatı veya tonu taklit etmesi istendiğinde.
\end{enumerate}

\vspace{4pt}
\begin{tcolorbox}[title=Cevap ve Açıklama, colback=gray!3!white, colframe=primarycolor!90!black, arc=2mm, fonttitle=\bfseries\sffamily, top=6pt, bottom=6pt]
    \textbf{Doğru Cevap: D}
    
    \vspace{2pt}
    Few-shot (Az Örnekli) istemleme, modele ne yapacağını sadece tarif etmek yerine girdilerin ve istenen çıktıların birkaç örneğini göstererek format ve ton uyumunu artırmak için kullanılır.
\end{tcolorbox}
\vspace{10pt}
% SORU 34
\needspace{5cm}
\vspace{8pt}
\noindent\textbf{\large\sffamily Soru 34.} \ Chain-of-Thought (CoT) tekniğinin 'Sihirli Kelimesi' veya ana mantığı nedir?

\begin{enumerate}[label=\Alph*), leftmargin=25pt, topsep=4pt, itemsep=2pt]
        \item Kısa ve öz cevap ver.
        \item Tablo formatında sun.
        \item Bir uzman gibi davran.
        \item Adım adım düşün.
\end{enumerate}

\vspace{4pt}
\begin{tcolorbox}[title=Cevap ve Açıklama, colback=gray!3!white, colframe=primarycolor!90!black, arc=2mm, fonttitle=\bfseries\sffamily, top=6pt, bottom=6pt]
    \textbf{Doğru Cevap: D}

    \vspace{2pt}
    Chain-of-Thought (CoT) tekniğinin temel mantığı, yapay zekanın karmaşık mantıksal veya matematiksel problemleri çözerken 'adım adım düşünmesini' sağlamaktır.
\end{tcolorbox}
\vspace{10pt}
% SORU 35
\needspace{5cm}
\vspace{8pt}
\noindent\textbf{\large\sffamily Soru 35.} \ Chunking (Parçalara Bölme) tekniği hangi problemi çözmek için kullanılır?

\begin{enumerate}[label=\Alph*), leftmargin=25pt, topsep=4pt, itemsep=2pt]
        \item Modelin eğitim verisini güncellemek.
        \item Modelin daha yaratıcı şiirler yazmasını sağlamak.
        \item Modelin halüsinasyon görmesini tamamen durdurmak.
        \item Darboğaz (Bottleneck) sorunu ve sistemin boğulmasını önlemek.
\end{enumerate}

\vspace{4pt}
\begin{tcolorbox}[title=Cevap ve Açıklama, colback=gray!3!white, colframe=primarycolor!90!black, arc=2mm, fonttitle=\bfseries\sffamily, top=6pt, bottom=6pt]
    \textbf{Doğru Cevap: D}

    \vspace{2pt}
    Chunking (Parçalara Bölme) tekniği, modelin tek seferde işleyemeyeceği büyüklükteki görevlerde sistemin boğulmasını ve sığ yanıtlar vermesini (darboğaz sorununu) önlemek amacıyla kullanılır.
\end{tcolorbox}
\vspace{10pt}
% SORU 36
\needspace{5cm}
\vspace{8pt}
\noindent\textbf{\large\sffamily Soru 36.} \ Kalite kontrol için kullanılan C.L.E.A.R. listesindeki 'Explicit' (Açık) kriteri neyi ifade eder?

\begin{enumerate}[label=\Alph*), leftmargin=25pt, topsep=4pt, itemsep=2pt]
        \item Kavramların doğru bir sırayla sunulması.
        \item Çıktı yönlendirmelerinin net olması (Örn: '50 kelimeyi geçme').
        \item Cevabın mantıklı olup olmadığının kontrol edilmesi.
        \item Gereksiz kelimelerden arındırılmış bir komut yazılması.
\end{enumerate}

\vspace{4pt}
\begin{tcolorbox}[title=Cevap ve Açıklama, colback=gray!3!white, colframe=primarycolor!90!black, arc=2mm, fonttitle=\bfseries\sffamily, top=6pt, bottom=6pt]
    \textbf{Doğru Cevap: B}

    \vspace{2pt}
    C.L.E.A.R. kontrol listesindeki 'Explicit' (Açık) kriteri, yapay zekaya verilen çıktı sınırlandırma komutlarının net ve ölçülebilir olmasını (örneğin '50 kelimeyi geçme') ifade eder.
\end{tcolorbox}
\vspace{10pt}
% SORU 37
\needspace{5cm}
\vspace{8pt}
\noindent\textbf{\large\sffamily Soru 37.} \ Görsel Üretimi (DALL-E, Midjourney vb.) için yapılan istemlerde hangi parametre 'Kritik' olarak belirtilmiştir?

\begin{enumerate}[label=\Alph*), leftmargin=25pt, topsep=4pt, itemsep=2pt]
        \item Zaman karmaşıklığını O(n log n) olarak optimize etmek.
        \item Kavramları değil, somut detayları betimlemek.
        \item Hedef kitleyi kesin olarak belirtmek.
        \item Kaynak ve hedef dilleri tanımlamak.
\end{enumerate}

\vspace{4pt}
\begin{tcolorbox}[title=Cevap ve Açıklama, colback=gray!3!white, colframe=primarycolor!90!black, arc=2mm, fonttitle=\bfseries\sffamily, top=6pt, bottom=6pt]
    \textbf{Doğru Cevap: B}

    \vspace{2pt}
    Görsel üretimi yapan yapay zeka modellerine komut yazılırken, soyut kavramlar yerine somut ve net detayların betimlenmesi çıktının kalitesini doğrudan artırır.
\end{tcolorbox}
\vspace{10pt}
% SORU 38
\needspace{5cm}
\vspace{8pt}
\noindent\textbf{\large\sffamily Soru 38.} \ İstem Mühendisliği 'Sentez Matrisi'ne göre CRAFTED ve CLEAR arasındaki ilişki nasıldır?

\begin{enumerate}[label=\Alph*), leftmargin=25pt, topsep=4pt, itemsep=2pt]
        \item CRAFTED ile inşa et, CLEAR ile test et ve iyileştir.
        \item CRAFTED ve CLEAR birbirine zıt iki farklı metoddur.
        \item CLEAR ile ilk taslağı oluştur, CRAFTED ile son formatı ver.
        \item CRAFTED sadece metin, CLEAR sadece görsel için kullanılır.
\end{enumerate}

\vspace{4pt}
\begin{tcolorbox}[title=Cevap ve Açıklama, colback=gray!3!white, colframe=primarycolor!90!black, arc=2mm, fonttitle=\bfseries\sffamily, top=6pt, bottom=6pt]
    \textbf{Doğru Cevap: A}

    \vspace{2pt}
    Sentez Matrisine göre, CRAFTED formülüyle istemi mükemmel şekilde inşa ederken, CLEAR listesiyle de üretilen çıktıyı test edip komutu rafine ederiz.
\end{tcolorbox}
\vspace{10pt}
% SORU 39
\needspace{5cm}
\vspace{8pt}
\noindent\textbf{\large\sffamily Soru 39.} \ Self-Consistency (Öz-Tutarlılık) tekniği en iyi hangi durumlarda sonuç verir?

\begin{enumerate}[label=\Alph*), leftmargin=25pt, topsep=4pt, itemsep=2pt]
        \item Modelin sistem komutlarını unutması durumunda.
        \item Kritik kararlar ve belirsizliğin olduğu matematik/mantık problemleri.
        \item Modelin bir çocuğa masal anlatması istendiğinde.
        \item Bir metni sadece başka bir dile çevirirken.
\end{enumerate}

\vspace{4pt}
\begin{tcolorbox}[title=Cevap ve Açıklama, colback=gray!3!white, colframe=primarycolor!90!black, arc=2mm, fonttitle=\bfseries\sffamily, top=6pt, bottom=6pt]
    \textbf{Doğru Cevap: B}

    \vspace{2pt}
    Self-Consistency (Öz-Tutarlılık) tekniği, özellikle kesin cevap gerektiren kritik kararlarda, matematik veya mantık problemlerinde en tutarlı sonucu seçmek için idealdir.
\end{tcolorbox}
\vspace{10pt}

\newpage
\section{Yapay Zeka Testi}
% SORU 40
\needspace{5cm}
\vspace{8pt}
\noindent\textbf{\large\sffamily Soru 40.} \ Üretken Yapay Zeka (ÜYZ) araçlarının akademik araştırmalardaki temel rolü aşağıdakilerden hangisi olarak tanımlanmıştır?

\begin{enumerate}[label=\Alph*), leftmargin=25pt, topsep=4pt, itemsep=2pt]
        \item Veri toplama sürecini insan müdahalesi olmadan tamamlayan bir otopilot.
        \item Bilimsel dürüstlük ilkelerinden muaf bir içerik üreticisi.
        \item Araştırma süreçlerinde verimlilik artırıcı sistematik bir destekçi.
        \item Tüm nihai kararları veren bağımsız bir araştırmacı.
\end{enumerate}

\vspace{4pt}
\begin{tcolorbox}[title=Cevap ve Açıklama, colback=gray!3!white, colframe=primarycolor!90!black, arc=2mm, fonttitle=\bfseries\sffamily, top=6pt, bottom=6pt]
    \textbf{Doğru Cevap: C}

    \vspace{2pt}
    Üretken Yapay Zeka (ÜYZ) araçları, akademik çalışmalarda nihai bir karar verici veya otopilot değil; araştırma süreçlerini hızlandıran ve verimliliği artıran sistematik bir destekçi (yardımcı pilot) rolündedir.
\end{tcolorbox}
\vspace{10pt}
% SORU 41
\needspace{5cm}
\vspace{8pt}
\noindent\textbf{\large\sffamily Soru 41.} \ ÜYZ modellerinin sunduğu çıktıların altında yatan 'Görünmez Önyargılar' (Bias) temel olarak nereden kaynaklanmaktadır?

\begin{enumerate}[label=\Alph*), leftmargin=25pt, topsep=4pt, itemsep=2pt]
        \item Kullanıcının yazdığı komutların (prompt) uzunluğundan.
        \item İnternet bağlantısındaki teknik aksaklıklardan.
        \item Modellerin eğitiminde kullanılan veri setlerinden.
        \item Kullanılan donanımın işlemci hızından.
\end{enumerate}

\vspace{4pt}
\begin{tcolorbox}[title=Cevap ve Açıklama, colback=gray!3!white, colframe=primarycolor!90!black, arc=2mm, fonttitle=\bfseries\sffamily, top=6pt, bottom=6pt]
    \textbf{Doğru Cevap: C}

    \vspace{2pt}
    Yapay zeka modellerinin sunduğu çıktılardaki görünmez önyargılar (bias), doğrudan modellerin eğitilmesinde kullanılan veri setlerindeki dengesizliklerden ve insan önyargılarından kaynaklanır.
\end{tcolorbox}
\vspace{10pt}
% SORU 42
\needspace{5cm}
\vspace{8pt}
\noindent\textbf{\large\sffamily Soru 42.} \ Etik ihlal durumunda araştırmacının 'Hata ÜYZ kaynaklıydı' savunmasının kabul edilmemesinin temel sebebi nedir?

\begin{enumerate}[label=\Alph*), leftmargin=25pt, topsep=4pt, itemsep=2pt]
        \item ÜYZ kullanımının tüm akademik dergiler tarafından yasaklanmış olması.
        \item Nihai sorumluluğun her zaman insan araştırmacıda olması.
        \item Yapay zeka araçlarının yasal bir kişiliğe sahip olması.
        \item ÜYZ sistemlerinin hata yapma ihtimalinin sıfır olması.
\end{enumerate}

\vspace{4pt}
\begin{tcolorbox}[title=Cevap ve Açıklama, colback=gray!3!white, colframe=primarycolor!90!black, arc=2mm, fonttitle=\bfseries\sffamily, top=6pt, bottom=6pt]
    \textbf{Doğru Cevap: B}

    \vspace{2pt}
    Yapay zeka araçları fikri ve entelektüel sorumluluk taşıyamayacakları için akademik çalışmalarda yazar olamazlar; dolayısıyla her türlü etik ve yasal sorumluluk nihai karar verici olan insana aittir.
\end{tcolorbox}
\vspace{10pt}
% SORU 43
\needspace{5cm}
\vspace{8pt}
\noindent\textbf{\large\sffamily Soru 43.} \ Şeffaflık (Transparency) ilkesi gereği, bir araştırmanın metodoloji bölümünde ÜYZ kullanımına dair hangi detayların verilmesi zorunludur?

\begin{enumerate}[label=\Alph*), leftmargin=25pt, topsep=4pt, itemsep=2pt]
        \item Aracın adı, sürüm bilgisi ve kullanım amacı.
        \item Yapay zekanın eğitiminde kullanılan tüm veri setlerinin listesi.
        \item Sadece yapay zekaya ödenen abonelik ücreti.
        \item Sadece yapay zeka ile yapılan sohbetlerin tam dökümü.
\end{enumerate}

\vspace{4pt}
\begin{tcolorbox}[title=Cevap ve Açıklama, colback=gray!3!white, colframe=primarycolor!90!black, arc=2mm, fonttitle=\bfseries\sffamily, top=6pt, bottom=6pt]
    \textbf{Doğru Cevap: A}

    \vspace{2pt}
    Şeffaflık (Transparency) ilkesi gereği, bilimsel çalışmalarda kullanılan yapay zeka aracının adı, sürüm/versiyon bilgisi ve tam olarak hangi amaçla kullanıldığı metodoloji bölümünde açıkça belirtilmelidir.
\end{tcolorbox}
\vspace{10pt}
% SORU 44
\needspace{5cm}
\vspace{8pt}
\noindent\textbf{\large\sffamily Soru 44.} \ Özen (Care) ilkesine göre, ÜYZ tarafından üretilen referanslar hakkında bilinen risk oranı nedir?

\begin{enumerate}[label=\Alph*), leftmargin=25pt, topsep=4pt, itemsep=2pt]
        \item \%50'ye varan oranda hatalı olabilir. (Not: Orijinal görselde metin "'a varan oranda hatalı olabilir" şeklindedir.)
        \item Hata riski sadece matematiksel formüllerde geçerlidir.
        \item hata payı ile kabul edilebilir.
        \item güvenilirdir.
\end{enumerate}

\vspace{4pt}
\begin{tcolorbox}[title=Cevap ve Açıklama, colback=gray!3!white, colframe=primarycolor!90!black, arc=2mm, fonttitle=\bfseries\sffamily, top=6pt, bottom=6pt]
    \textbf{Doğru Cevap: A}

    \vspace{2pt}
    Özen (Care) ilkesine göre, yapay zeka modelleri tarafından uydurulan veya akademik kaynakçaların hatalı olma riski yaklaşık %30'a ulaşabildiğinden, her referansın doğruluğu manuel olarak kontrol edilmelidir.
\end{tcolorbox}
\vspace{10pt}
% SORU 45
\needspace{5cm}
\vspace{8pt}
\noindent\textbf{\large\sffamily Soru 45.} \ ÜYZ araçlarının akademik çalışmalarda 'yazar' olarak kaydedilememesinin temel nedeni nedir?

\begin{enumerate}[label=\Alph*), leftmargin=25pt, topsep=4pt, itemsep=2pt]
        \item Üretilen metinlerin her zaman çok kısa olması.
        \item Yeterli kelime haznesine sahip olmamaları.
        \item Ortak yazarlık kriterlerini ve entelektüel sorumluluğu karşılayamamaları.
        \item Türk Dil Kurumu sözlüğünde 'yazar' tanımına uymamaları.
\end{enumerate}

\vspace{4pt}
\begin{tcolorbox}[title=Cevap ve Açıklama, colback=gray!3!white, colframe=primarycolor!90!black, arc=2mm, fonttitle=\bfseries\sffamily, top=6pt, bottom=6pt]
    \textbf{Doğru Cevap: C}

    \vspace{2pt}
    Yapay zeka araçları, ortak yazarlık kriterlerini (entelektüel katkı ve sorumluluk alma) karşılayamadıkları için bilimsel yayınlarda yazar (co-author) olarak listelenemezler.
\end{tcolorbox}
\vspace{10pt}
% SORU 46
\needspace{5cm}
\vspace{8pt}
\noindent\textbf{\large\sffamily Soru 46.} \ Gizlilik ve Mahremiyet ilkesi çerçevesinde, hassas verilerin ÜYZ modellerine yüklenmeden önce geçmesi gereken işlem nedir?

\begin{enumerate}[label=\Alph*), leftmargin=25pt, topsep=4pt, itemsep=2pt]
        \item Verilerin İngilizceye çevrilmesi.
        \item Verilerin sadece gece saatlerinde yüklenmesi.
        \item Verilerin bir Word dosyası olarak yüklenmesi.
        \item Anonimleştirme (Kişisel verilerin temizlenmesi).
\end{enumerate}

\vspace{4pt}
\begin{tcolorbox}[title=Cevap ve Açıklama, colback=gray!3!white, colframe=primarycolor!90!black, arc=2mm, fonttitle=\bfseries\sffamily, top=6pt, bottom=6pt]
    \textbf{Doğru Cevap: D}

    \vspace{2pt}
    Gizlilik ve Mahremiyet ilkesi çerçevesinde, kişisel veya hassas verilerin yapay zeka modellerine yüklenip işlenmesinden önce mutlaka anonimleştirilmesi (kimlik bilgilerinden arındırılması) gerekir.
\end{tcolorbox}
\vspace{10pt}
% SORU 47
\needspace{5cm}
\vspace{8pt}
\noindent\textbf{\large\sffamily Soru 47.} \ 'Dengeyi Kurmak: Pilot vs. Otopilot' görselinde vurgulanan ana mesaj nedir?

\begin{enumerate}[label=\Alph*), leftmargin=25pt, topsep=4pt, itemsep=2pt]
        \item Otopilot modu bilimsel araştırmalarda en güvenli yöntemdir.
        \item Yapay zeka geliştikçe insan araştırmacıya gerek kalmayacaktır.
        \item İnsan araştırmacı her zaman kaptandır, ÜYZ ise yardımcı pilottur.
        \item Araştırmacılar sadece veri girişinden sorumludur.
\end{enumerate}

\vspace{4pt}
\begin{tcolorbox}[title=Cevap ve Açıklama, colback=gray!3!white, colframe=primarycolor!90!black, arc=2mm, fonttitle=\bfseries\sffamily, top=6pt, bottom=6pt]
    \textbf{Doğru Cevap: C}

    \vspace{2pt}
    Akademik çalışmalarda insan araştırmacı her zaman süreci yöneten ve kararları veren 'kaptan/pilot' konumundadır; yapay zeka ise sadece ona yardımcı olan bir 'otopilot/yardımcı pilot' rolündedir.
\end{tcolorbox}
\vspace{10pt}
% SORU 48
\needspace{5cm}
\vspace{8pt}
\noindent\textbf{\large\sffamily Soru 48.} \ Sorumlu Kullanım Kontrol Listesi'ne (Checklist) göre, yayını göndermeden önce yapılması gereken son adımlardan biri hangisidir?

\begin{enumerate}[label=\Alph*), leftmargin=25pt, topsep=4pt, itemsep=2pt]
        \item Tüm metni yapay zekaya yeniden yazdırtmak.
        \item Kaynakça listesindeki her bir eserin gerçek ve doğru olduğunu teyit etmek.
        \item Yapay zeka çalışmayı beğenip beğenmediğini sormak.
        \item Yapay zekanın kullandığı elektrik miktarını hesaplamak.
\end{enumerate}

\vspace{4pt}
\begin{tcolorbox}[title=Cevap ve Açıklama, colback=gray!3!white, colframe=primarycolor!90!black, arc=2mm, fonttitle=\bfseries\sffamily, top=6pt, bottom=6pt]
    \textbf{Doğru Cevap: B}

    \vspace{2pt}
    Sorumlu kullanım kontrol listesine göre, akademik yayını göndermeden önce kaynakçada yer alan tüm referansların uydurma olmadığını ve gerçek eserlere dayandığını manuel olarak doğrulamak kritik bir son adımdır.
\end{tcolorbox}
\vspace{10pt}
% SORU 49
\needspace{5cm}
\vspace{8pt}
\noindent\textbf{\large\sffamily Soru 49.} \ Adalet ve Saygı (Justice \& Respect) sütununda belirtilen 'Disiplin' gereği aşağıdakilerden hangisi bir akademik suç olarak tanımlanmıştır?

\begin{enumerate}[label=\Alph*), leftmargin=25pt, topsep=4pt, itemsep=2pt]
        \item Birden fazla ÜYZ aracını aynı anda kullanmak.
        \item ÜYZ araçlarının ücretsiz versiyonlarını kullanmak.
        \item ÜYZ kullanarak sadece dil düzeltmesi yapmak.
        \item ÜYZ kullanımının gizlenmesi.
\end{enumerate}

\vspace{4pt}
\begin{tcolorbox}[title=Cevap ve Açıklama, colback=gray!3!white, colframe=primarycolor!90!black, arc=2mm, fonttitle=\bfseries\sffamily, top=6pt, bottom=6pt]
    \textbf{Doğru Cevap: D}

    \vspace{2pt}
    Yapay zeka kullanıldığını saklamak veya gizlemek ciddi bir akademik etik ihlali ve suç olarak kabul edilir.
\end{tcolorbox}
\vspace{10pt}
% SORU 50
\needspace{5cm}
\vspace{8pt}
\noindent\textbf{\large\sffamily Soru 50.} \ Görsel yapay zekanın temel çalışma prensibi aşağıdakilerden hangisidir?

\begin{enumerate}[label=\Alph*), leftmargin=25pt, topsep=4pt, itemsep=2pt]
        \item Metin komutlarından özgün görseller yaratmak
        \item Sadece var olan fotoğrafları düzenlemek
        \item Elle çizilen resimleri dijital ortama aktarmak
        \item İnternetteki görselleri telif haklarına bakmadan kopyalamak
\end{enumerate}

\vspace{4pt}
\begin{tcolorbox}[title=Cevap ve Açıklama, colback=gray!3!white, colframe=primarycolor!90!black, arc=2mm, fonttitle=\bfseries\sffamily, top=6pt, bottom=6pt]
    \textbf{Doğru Cevap: A}

    \vspace{2pt}
    Görsel yapay zekanın (DALL-E, Midjourney, Stable Diffusion vb.) temel çalışma prensibi, kullanıcının girdiği metinsel tanımlardan (prompt) yola çıkarak tamamen özgün ve yeni görseller yaratmaktır.
\end{tcolorbox}
\vspace{10pt}
% SORU 51
\needspace{5cm}
\vspace{8pt}
\noindent\textbf{\large\sffamily Soru 51.} \ DALL-E 2 modelinin kullanıcılar için sunduğu en belirgin avantaj hangisidir?

\begin{enumerate}[label=\Alph*), leftmargin=25pt, topsep=4pt, itemsep=2pt]
        \item Yalnızca topluluk üzerinden erişilebilir olması
        \item Sadece profesyonel yazılımcılara açık olması
        \item Görsellerin sadece siyah-beyaz üretilmesi
        \item Teknik bilgi gerektirmeyen kullanıcı dostu arayüz
\end{enumerate}

\vspace{4pt}
\begin{tcolorbox}[title=Cevap ve Açıklama, colback=gray!3!white, colframe=primarycolor!90!black, arc=2mm, fonttitle=\bfseries\sffamily, top=6pt, bottom=6pt]
    \textbf{Doğru Cevap: D}

    \vspace{2pt}
    DALL-E 2 modeli, teknik veya yazılımsal bilgi gerektirmeyen son derece basit, web tabanlı ve kullanıcı dostu bir arayüz sunduğu için geniş kitleler tarafından tercih edilmiştir.
\end{tcolorbox}
\vspace{10pt}
% SORU 52
\needspace{5cm}
\vspace{8pt}
\noindent\textbf{\large\sffamily Soru 52.} \ Sanatsal ve sinematik estetiği ile ön plana çıkan, topluluk odaklı platform aşağıdakilerden hangisidir?

\begin{enumerate}[label=\Alph*), leftmargin=25pt, topsep=4pt, itemsep=2pt]
        \item Midjourney
        \item DALL-E 2
        \item Adobe Firefly
        \item Stable Diffusion
\end{enumerate}

\vspace{4pt}
\begin{tcolorbox}[title=Cevap ve Açıklama, colback=gray!3!white, colframe=primarycolor!90!black, arc=2mm, fonttitle=\bfseries\sffamily, top=6pt, bottom=6pt]
    \textbf{Doğru Cevap: A}

    \vspace{2pt}
    Midjourney, ürettiği görsellerdeki yüksek sanatsal detay, sinematik estetik ve Discord tabanlı topluluk yapısıyla görsel yapay zeka araçları arasında öne çıkmaktadır.
\end{tcolorbox}
\vspace{10pt}
% SORU 53
\needspace{5cm}
\vspace{8pt}
\noindent\textbf{\large\sffamily Soru 53.} \ Stable Diffusion modelini diğer yapay zeka araçlarından ayıran temel teknik özellik nedir?

\begin{enumerate}[label=\Alph*), leftmargin=25pt, topsep=4pt, itemsep=2pt]
        \item Ücretli abonelik zorunluluğu olan tek model olması
        \item Açık kaynak kodlu ve özelleştirilebilir olması
        \item Sadece mobil cihazlarda çalışabilmesi
        \item Telif hakkı riskini tamamen ortadan kaldırması
\end{enumerate}

\vspace{4pt}
\begin{tcolorbox}[title=Cevap ve Açıklama, colback=gray!3!white, colframe=primarycolor!90!black, arc=2mm, fonttitle=\bfseries\sffamily, top=6pt, bottom=6pt]
    \textbf{Doğru Cevap: B}

    \vspace{2pt}
    Stable Diffusion modelini diğer ticari araçlardan ayıran en temel fark, açık kaynak kodlu olması ve kullanıcıların kendi yerel donanımlarında çalıştırıp özelleştirebilmelerine olanak tanımasıdır.
\end{tcolorbox}
\vspace{10pt}
% SORU 54
\needspace{5cm}
\vspace{8pt}
\noindent\textbf{\large\sffamily Soru 54.} \ Adobe Firefly'ın özellikle profesyonel tasarımcılar tarafından tercih edilme sebebi nedir?

\begin{enumerate}[label=\Alph*), leftmargin=25pt, topsep=4pt, itemsep=2pt]
        \item Telif hakkı riski olmayan içerikler sunması
        \item Sadece video düzenleme için kullanılması
        \item Tamamen ücretsiz ve sınırsız olması
        \item İnternet bağlantısı olmadan çalışması
\end{enumerate}

\vspace{4pt}
\begin{tcolorbox}[title=Cevap ve Açıklama, colback=gray!3!white, colframe=primarycolor!90!black, arc=2mm, fonttitle=\bfseries\sffamily, top=6pt, bottom=6pt]
    \textbf{Doğru Cevap: A}

    \vspace{2pt}
    Adobe Firefly, lisanslı ve telif hakkı temiz veri setleriyle eğitildiği için profesyonel tasarımcılara ve kurumlara telif hakkı riski taşımayan güvenli ticari görseller üretme avantajı sağlar.
\end{tcolorbox}
\vspace{10pt}
% SORU 55
\needspace{5cm}
\vspace{8pt}
\noindent\textbf{\large\sffamily Soru 55.} \ Sunumda 'Yaratıcı Deneyim' başlığı altında hangi araçlar bir arada anılmıştır?

\begin{enumerate}[label=\Alph*), leftmargin=25pt, topsep=4pt, itemsep=2pt]
        \item Photoshop ve Illustrator
        \item GauGAN ve NightCafe
        \item DALL-E 2 ve Midjourney
        \item Leonardo AI ve Artbreeder
\end{enumerate}

\vspace{4pt}
\begin{tcolorbox}[title=Cevap ve Açıklama, colback=gray!3!white, colframe=primarycolor!90!black, arc=2mm, fonttitle=\bfseries\sffamily, top=6pt, bottom=6pt]
    \textbf{Doğru Cevap: D}

    \vspace{2pt}
    Yaratıcı Deneyim başlığı altında GauGAN ve NightCafe gibi araçlar yer alırken, sunumda bu başlık altında Leonardo AI ve Artbreeder araçları da bir arada anılmıştır. Doğru cevap D şıkkıdır.
\end{tcolorbox}
\vspace{10pt}
% SORU 56
\needspace{5cm}
\vspace{8pt}
\noindent\textbf{\large\sffamily Soru 56.} \ Bir yapay zeka aracı seçilirken dikkat edilmesi gereken kriterler arasında aşağıdakilerden hangisi yer almaz?

\begin{enumerate}[label=\Alph*), leftmargin=25pt, topsep=4pt, itemsep=2pt]
        \item Kullanıcının eğitim seviyesi
        \item Telif hakları ve güvenlik standartları
        \item Maliyet ve kullanım hızı
        \item Çıktı kalitesi ve stil uygunluğu
\end{enumerate}

\vspace{4pt}
\begin{tcolorbox}[title=Cevap ve Açıklama, colback=gray!3!white, colframe=primarycolor!90!black, arc=2mm, fonttitle=\bfseries\sffamily, top=6pt, bottom=6pt]
    \textbf{Doğru Cevap: A}

    \vspace{2pt}
    Bir yapay zeka aracı seçilirken kullanıcının kendi eğitim seviyesi teknik bir kriter değildir; asıl olarak maliyet, hız, telif hakları, çıktı kalitesi ve stil uygunluğu gibi kriterlere bakılır.
\end{tcolorbox}
\vspace{10pt}
% SORU 57
\needspace{5cm}
\vspace{8pt}
\noindent\textbf{\large\sffamily Soru 57.} \ Görsel üretim dünyasının geleceği hakkında yapılan öngörü aşağıdakilerden hangisidir?

\begin{enumerate}[label=\Alph*), leftmargin=25pt, topsep=4pt, itemsep=2pt]
        \item Geleneksel sanatın tamamen sona ermesi
        \item Tüm görsel üretimin tek bir şirket tarafından kontrol edilmesi
        \item İnsanların artık tasarım yapmasına gerek kalmaması
        \item İnsan ve yapay zeka iş birliği
\end{enumerate}

\vspace{4pt}
\begin{tcolorbox}[title=Cevap ve Açıklama, colback=gray!3!white, colframe=primarycolor!90!black, arc=2mm, fonttitle=\bfseries\sffamily, top=6pt, bottom=6pt]
    \textbf{Doğru Cevap: D}

    \vspace{2pt}
    Görsel yapay zeka teknolojilerinin geleceğinde, geleneksel sanatın tamamen yok olması değil; insan yaratıcılığı ile yapay zekanın gücünün birleştiği bir insan-makine iş birliği öngörülmektedir.
\end{tcolorbox}
\vspace{10pt}
% SORU 58
\needspace{5cm}
\vspace{8pt}
\noindent\textbf{\large\sffamily Soru 58.} \ Betül Hiçdurmaz tarafından hazırlanan bu sunumda vurgulanan temel tema nedir?

\begin{enumerate}[label=\Alph*), leftmargin=25pt, topsep=4pt, itemsep=2pt]
        \item Sosyal medya kullanımının yaygınlaşması
        \item Görsel üretiminde yapay zeka devrimi
        \item Veri tabanı yönetimi yöntemleri
        \item Eski resimlerin restorasyonu
\end{enumerate}

\vspace{4pt}
\begin{tcolorbox}[title=Cevap ve Açıklama, colback=gray!3!white, colframe=primarycolor!90!black, arc=2mm, fonttitle=\bfseries\sffamily, top=6pt, bottom=6pt]
    \textbf{Doğru Cevap: B}

    \vspace{2pt}
    Betül Hiçdurmaz tarafından hazırlanan sunumda, yapay zekanın görsel üretim alanında gerçekleştirdiği devrim, kullanılan araçlar ve bunların tasarım süreçlerine etkileri ele alınmıştır.
\end{tcolorbox}
\vspace{10pt}
% SORU 59
\needspace{5cm}
\vspace{8pt}
\noindent\textbf{\large\sffamily Soru 59.} \ GauGAN ve NightCafe gibi araçlar hangi başlık altında sınıflandırılmıştır?

\begin{enumerate}[label=\Alph*), leftmargin=25pt, topsep=4pt, itemsep=2pt]
        \item Yüksek Kaliteli Ana Modeller
        \item Telif Hakkı Garanti Eden Araçlar
        \item Ücretli Tasarım Yazılımları
        \item Özel Kullanım Alanları
\end{enumerate}

\vspace{4pt}
\begin{tcolorbox}[title=Cevap ve Açıklama, colback=gray!3!white, colframe=primarycolor!90!black, arc=2mm, fonttitle=\bfseries\sffamily, top=6pt, bottom=6pt]
    \textbf{Doğru Cevap: D}

    \vspace{2pt}
    GauGAN (karalamaları gerçekçi manzaralara dönüştürme) ve NightCafe (farklı sanatsal stiller uygulama) gibi araçlar, genel amaçlı modellerden ziyade belirli özel kullanım alanlarına hitap eder.
\end{tcolorbox}
\vspace{10pt}
% SORU 60
\needspace{5cm}
\vspace{8pt}
\noindent\textbf{\large\sffamily Soru 60.} \ Araştırmada Üretken Yapay Zeka (ÜYZ) kullanımında 'Şeffaflık' ilkesi gereği, metodoloji bölümüne hangi detayların eklenmesi zorunludur?

\begin{enumerate}[label=\Alph*), leftmargin=25pt, topsep=4pt, itemsep=2pt]
        \item Yapay zekaya verilen tüm komutların (prompt) tam listesi.
        \item Sadece aracın üretici firmasının adı ve genel merkezi.
        \item Aracın eğitiminde kullanılan veri setlerinin boyutu.
        \item Kullanılan aracın adı, sürüm bilgisi ve kullanım amacı.
\end{enumerate}

\vspace{4pt}
\begin{tcolorbox}[title=Cevap ve Açıklama, colback=gray!3!white, colframe=primarycolor!90!black, arc=2mm, fonttitle=\bfseries\sffamily, top=6pt, bottom=6pt]
    \textbf{Doğru Cevap: D}

    \vspace{2pt}
    Araştırmalarda Üretken Yapay Zeka kullanımında şeffaflık ilkesi gereği, metodoloji bölümünde kullanılan yapay zeka aracının adı, sürümü (versiyonu) ve kullanım amacı açıkça belirtilmelidir.
\end{tcolorbox}
\vspace{10pt}
% SORU 61
\needspace{5cm}
\vspace{8pt}
\noindent\textbf{\large\sffamily Soru 61.} \ ÜYZ tarafından üretilen bir içerikte bilimsel hata veya intihal tespit edilmesi durumunda sorumluluk kime aittir?

\begin{enumerate}[label=\Alph*), leftmargin=25pt, topsep=4pt, itemsep=2pt]
        \item ÜYZ modelinin eğitiminde kullanılan veri sağlayıcılara.
        \item Üniversitenin veya kurumun etik kuruluna.
        \item Yazılımı geliştiren teknoloji şirketine.
        \item Araştırmanın kendisine (nihai karar verici olan insana).
\end{enumerate}

\vspace{4pt}
\begin{tcolorbox}[title=Cevap ve Açıklama, colback=gray!3!white, colframe=primarycolor!90!black, arc=2mm, fonttitle=\bfseries\sffamily, top=6pt, bottom=6pt]
    \textbf{Doğru Cevap: D}

    \vspace{2pt}
    Yapay zeka entelektüel bir sorumluluk üstlenemeyeceği için, üretilen çıktılarda oluşabilecek bilimsel hata, intihal veya yanlış beyanlardan doğrudan çalışmayı yürüten insan araştırmacı sorumludur.
\end{tcolorbox}
\vspace{10pt}
% SORU 62
\needspace{5cm}
\vspace{8pt}
\noindent\textbf{\large\sffamily Soru 62.} \ Araştırma verilerine göre ÜYZ sistemlerinin ürettiği referansların hatalı olma riski yaklaşık yüzde kaçtır?

\begin{enumerate}[label=\Alph*), leftmargin=25pt, topsep=4pt, itemsep=2pt]
        \item 50\%
        \item 30\%
        \item 75\%
        \item 10\%
\end{enumerate}

\vspace{4pt}
\begin{tcolorbox}[title=Cevap ve Açıklama, colback=gray!3!white, colframe=primarycolor!90!black, arc=2mm, fonttitle=\bfseries\sffamily, top=6pt, bottom=6pt]
    \textbf{Doğru Cevap: B}

    \vspace{2pt}
    Araştırmalara göre, üretken yapay zeka modelleri tarafından uydurulan veya hatalı üretilen akademik kaynakça ve referansların oranı yaklaşık %30'dur.
\end{tcolorbox}
\vspace{10pt}
% SORU 63
\needspace{5cm}
\vspace{8pt}
\noindent\textbf{\large\sffamily Soru 63.} \ Adobe Firefly'ı profesyonel tasarımcılar için diğer görsel ÜYZ araçlarından ayıran temel avantaj nedir?

\begin{enumerate}[label=\Alph*), leftmargin=25pt, topsep=4pt, itemsep=2pt]
        \item Telif hakkı riski olmayan içerikler sunması ve kurumsal entegrasyonu.
        \item Sadece metinden video üretmeye odaklanması.
        \item Sadece soyut ve sanatsal görseller üretmesi.
        \item Açık kaynak kodlu ve tamamen ücretsiz olması.
\end{enumerate}

\vspace{4pt}
\begin{tcolorbox}[title=Cevap ve Açıklama, colback=gray!3!white, colframe=primarycolor!90!black, arc=2mm, fonttitle=\bfseries\sffamily, top=6pt, bottom=6pt]
    \textbf{Doğru Cevap: A}

    \vspace{2pt}
    Adobe Firefly, Adobe Stock ve telifsiz görsellerle eğitildiğinden, profesyonel tasarımcılara kurumsal entegrasyon kolaylığı ve telif hakkı riski olmayan güvenli bir çalışma ortamı sunar.
\end{tcolorbox}
\vspace{10pt}
% SORU 64
\needspace{5cm}
\vspace{8pt}
\noindent\textbf{\large\sffamily Soru 64.} \ Özel bir 'Gem' inşa ederken sistemin 'Taşıyıcı Kolonları' olarak adlandırılan 4 temel yapı taşı hangileridir?

\begin{enumerate}[label=\Alph*), leftmargin=25pt, topsep=4pt, itemsep=2pt]
        \item Analiz, Tasarım, Geliştirme, Test.
        \item Hız, Maliyet, Güvenlik, Erişilebilirlik.
        \item Veri, Donanım, Yazılım, Kullanıcı.
        \item Karakter, Görev, Bağlam, Biçim.
\end{enumerate}

\vspace{4pt}
\begin{tcolorbox}[title=Cevap ve Açıklama, colback=gray!3!white, colframe=primarycolor!90!black, arc=2mm, fonttitle=\bfseries\sffamily, top=6pt, bottom=6pt]
    \textbf{Doğru Cevap: D}

    \vspace{2pt}
    Özel bir Gem asistanı tasarlarken sistemi yapılandıran ve promptun temelini oluşturan 4 kolon: Karakter (Persona), Görev (Task), Bağlam (Context) ve Biçim (Format) olarak tanımlanır.
\end{tcolorbox}
\vspace{10pt}
% SORU 65
\needspace{5cm}
\vspace{8pt}
\noindent\textbf{\large\sffamily Soru 65.} \ İstem mühendisliğinde 'Chain of Thought' (Düşünce Zinciri) yönteminin temel amacı nedir?

\begin{enumerate}[label=\Alph*), leftmargin=25pt, topsep=4pt, itemsep=2pt]
        \item Çıktının kelime sayısını artırarak daha uzun metinler elde etmek.
        \item Modelin adım adım akıl yürüterek karmaşık problemlere daha doğru yanıtlar vermesini sağlamak.
        \item Sadece dil çevirisi görevlerinde başarıyı artırmak.
        \item Modelin çıktıyı daha hızlı üretmesini sağlamak.
\end{enumerate}

\vspace{4pt}
\begin{tcolorbox}[title=Cevap ve Açıklama, colback=gray!3!white, colframe=primarycolor!90!black, arc=2mm, fonttitle=\bfseries\sffamily, top=6pt, bottom=6pt]
    \textbf{Doğru Cevap: B}

    \vspace{2pt}
    Düşünce Zinciri (Chain of Thought) yöntemi, modelin bir problemi çözerken ara adımları ve mantık silsilesini takip etmesini sağlayarak karmaşık problemlerde daha doğru yanıtlar üretmesini amaçlar.
\end{tcolorbox}
\vspace{10pt}
% SORU 66
\needspace{5cm}
\vspace{8pt}
\noindent\textbf{\large\sffamily Soru 66.} \ Üretken Çekişmeli Ağlar (GAN) mimarisinde birbiriyle rekabet eden iki ana bileşen hangileridir?

\begin{enumerate}[label=\Alph*), leftmargin=25pt, topsep=4pt, itemsep=2pt]
        \item Oluşturucu (Generator) ve Ayırt Edici (Discriminator).
        \item Gürültü Ekleyici ve Gürültü Çıkarıcı.
        \item Sorgu (Query) ve Anahtar (Key).
        \item Veri Temizleyici ve Veri Düzenleyici.
\end{enumerate}

\vspace{4pt}
\begin{tcolorbox}[title=Cevap ve Açıklama, colback=gray!3!white, colframe=primarycolor!90!black, arc=2mm, fonttitle=\bfseries\sffamily, top=6pt, bottom=6pt]
    \textbf{Doğru Cevap: A}

    \vspace{2pt}
    GAN (Üretken Çekişmeli Ağlar) mimarisinde, yeni veri üreten Oluşturucu (Generator) ile üretilen verinin gerçek mi yoksa sahte mi olduğunu denetleyen Ayırt Edici (Discriminator) olmak üzere iki sinir ağı birbiriyle yarışır.
\end{tcolorbox}
\vspace{10pt}
% SORU 67
\needspace{5cm}
\vspace{8pt}
\noindent\textbf{\large\sffamily Soru 67.} \ Transformer mimarisindeki 'Self-Attention' mekanizması için hangisi söylenemez?

\begin{enumerate}[label=\Alph*), leftmargin=25pt, topsep=4pt, itemsep=2pt]
        \item Girdi dizisindeki farklı öğelerin birbirine göre önemini hesaplar.
        \item Q, K ve V olmak üzere üç ana vektör bileşeni kullanır.
        \item Uzun mesafeli bağlamları yakalamada etkilidir.
        \item Paralel hesaplamaya uygun değildir ve RNN'lerden daha yavaştır.
\end{enumerate}

\vspace{4pt}
\begin{tcolorbox}[title=Cevap ve Açıklama, colback=gray!3!white, colframe=primarycolor!90!black, arc=2mm, fonttitle=\bfseries\sffamily, top=6pt, bottom=6pt]
    \textbf{Doğru Cevap: D}

    \vspace{2pt}
    Transformer mimarisinin RNN ve LSTM gibi eski yapılara kıyasla en büyük avantajlarından biri, dizi öğeleri arasındaki ilişkileri paralel olarak işleyebilmesidir.
\end{tcolorbox}
\vspace{10pt}
% SORU 68
\needspace{5cm}
\vspace{8pt}
\noindent\textbf{\large\sffamily Soru 68.} \ Üretken yapay zekanın tarihsel gelişiminde 'Transformer' mimarisinin tanıtıldığı ve modern AI dönemini başlatan yıl hangisidir?

\begin{enumerate}[label=\Alph*), leftmargin=25pt, topsep=4pt, itemsep=2pt]
        \item 1997
        \item 2022
        \item 2014
        \item 2017
\end{enumerate}

\vspace{4pt}
\begin{tcolorbox}[title=Cevap ve Açıklama, colback=gray!3!white, colframe=primarycolor!90!black, arc=2mm, fonttitle=\bfseries\sffamily, top=6pt, bottom=6pt]
    \textbf{Doğru Cevap: D}

    \vspace{2pt}
    Üretken yapay zekada bir dönüm noktası kabul edilen ve self-attention mekanizmasını tanıtan 'Attention Is All You Need' makalesiyle Transformer mimarisi 2017 yılında yayımlanmıştır.
\end{tcolorbox}
\vspace{10pt}
% SORU 69
\needspace{5cm}
\vspace{8pt}
\noindent\textbf{\large\sffamily Soru 69.} \ Yapay zeka etiğindeki 'Gizlilik ve Mahremiyet' sütunu kapsamında, katılımcı verileri ÜYZ sistemine yüklenmeden önce hangi işlem mutlaka yapılmalıdır?

\begin{enumerate}[label=\Alph*), leftmargin=25pt, topsep=4pt, itemsep=2pt]
        \item Veriler alfabetik olarak sıralanmalıdır.
        \item Verilerin yedeği alınmalıdır.
        \item Veriler tamamen anonimleştirilmelidir.
        \item Veriler sadece İngilizceye çevrilmelidir.
\end{enumerate}

\vspace{4pt}
\begin{tcolorbox}[title=Cevap ve Açıklama, colback=gray!3!white, colframe=primarycolor!90!black, arc=2mm, fonttitle=\bfseries\sffamily, top=6pt, bottom=6pt]
    \textbf{Doğru Cevap: C}

    \vspace{2pt}
    Yapay zeka etiği gereği, hassas veya kişisel verilerin sisteme yüklenmeden önce kimlik bilgilerinden arındırılması yani tamamen anonimleştirilmesi zorunludur.
\end{tcolorbox}
\vspace{10pt}
% SORU 70
\needspace{5cm}
\vspace{8pt}
\noindent\textbf{\large\sffamily Soru 70.} \ Ayrıştırıcı (Discriminative) modeller ile Üretken (Generative) modeller arasındaki temel fark aşağıdakilerden hangisidir?

\begin{enumerate}[label=\Alph*), leftmargin=25pt, topsep=4pt, itemsep=2pt]
        \item Ayrıştırıcı modeller veriyi sınıflandırırken, üretken modeller eğitildikleri veri kümesine dayanarak tamamen yeni içerikler oluşturur.
        \item Ayrıştırıcı modeller daha güçlü donanım gerektirirken, üretken modeller standart bilgisayarlarda çalışabilir.
        \item Üretken modeller sadece metin tabanlı çalışırken, ayrıştırıcı modeller tüm medya türlerini işleyebilir.
        \item Üretken modeller açık talimatlar (programlama) ile çalışırken, ayrıştırıcı modeller kendi kendine öğrenir.
\end{enumerate}

\vspace{4pt}
\begin{tcolorbox}[title=Cevap ve Açıklama, colback=gray!3!white, colframe=primarycolor!90!black, arc=2mm, fonttitle=\bfseries\sffamily, top=6pt, bottom=6pt]
    \textbf{Doğru Cevap: A}

    \vspace{2pt}
    Ayrıştırıcı (Discriminative) modeller verinin hangi sınıfa ait olduğunu (sınıf sınırlarını) öğrenirken; Üretken (Generative) modeller verinin dağılımını öğrenerek tamamen yeni ve özgün içerikler oluşturur.
\end{tcolorbox}
\vspace{10pt}
% SORU 71
\needspace{5cm}
\vspace{8pt}
\noindent\textbf{\large\sffamily Soru 71.} \ DeepSeek R1 modelinin matematiksel problemlerde o3 gibi rakiplerine karşı sunduğu en belirgin kullanıcı deneyimi avantajı nedir?

\begin{enumerate}[label=\Alph*), leftmargin=25pt, topsep=4pt, itemsep=2pt]
        \item Matematik yarışmalarında her zaman en yüksek puanı alması.
        \item Çözüm adımlarını ve düşünme sürecini şeffaf bir şekilde göstermesi.
        \item Yanıtları rakiplerinden çok daha hızlı ve kısa bir şekilde üretmesi.
        \item En pahalı abonelik modeline sahip olması.
\end{enumerate}

\vspace{4pt}
\begin{tcolorbox}[title=Cevap ve Açıklama, colback=gray!3!white, colframe=primarycolor!90!black, arc=2mm, fonttitle=\bfseries\sffamily, top=6pt, bottom=6pt]
    \textbf{Doğru Cevap: B}

    \vspace{2pt}
    DeepSeek R1 modeli, akıl yürütme (reasoning) adımlarını ve düşünme sürecini (thought process) şeffaf bir şekilde arayüzde göstererek kullanıcılara açıklanabilir bir deneyim sunar.
\end{tcolorbox}
\vspace{10pt}
% SORU 72
\needspace{5cm}
\vspace{8pt}
\noindent\textbf{\large\sffamily Soru 72.} \ Transformer mimarisinin 'Öz-dikkat' (Self-attention) mekanizmasında kullanılan üç ana bileşen hangisidir?

\begin{enumerate}[label=\Alph*), leftmargin=25pt, topsep=4pt, itemsep=2pt]
        \item Generator, Discriminator, Noise
        \item Encoder, Decoder, Latent
        \item Query (Q), Key (K), Value (V)
        \item Input, Hidden, Output
\end{enumerate}

\vspace{4pt}
\begin{tcolorbox}[title=Cevap ve Açıklama, colback=gray!3!white, colframe=primarycolor!90!black, arc=2mm, fonttitle=\bfseries\sffamily, top=6pt, bottom=6pt]
    \textbf{Doğru Cevap: C}

    \vspace{2pt}
    Transformer mimarisindeki Öz-Dikkat (Self-Attention) mekanizması, girdi kelimelerinin birbirleriyle olan anlamsal ilişkilerini Query (Sorgu), Key (Anahtar) ve Value (Değer) vektörleri aracılığıyla hesaplar.
\end{tcolorbox}
\vspace{10pt}
% SORU 73
\needspace{5cm}
\vspace{8pt}
\noindent\textbf{\large\sffamily Soru 73.} \ 2025 yılı yapay zeka model karşılaştırmalarına göre, Claude Sonnet 4.5 hangi alandaki uzmanlığı ile öne çıkmaktadır?

\begin{enumerate}[label=\Alph*), leftmargin=25pt, topsep=4pt, itemsep=2pt]
        \item Microsoft Office uygulamalarıyla tam entegrasyon
        \item Kodlama, ajanlar ve bilgisayar kullanımı
        \item En düşük maliyetli API kullanımı
        \item Genel amaçlı muhakeme ve bilgi işi
\end{enumerate}

\vspace{4pt}
\begin{tcolorbox}[title=Cevap ve Açıklama, colback=gray!3!white, colframe=primarycolor!90!black, arc=2mm, fonttitle=\bfseries\sffamily, top=6pt, bottom=6pt]
    \textbf{Doğru Cevap: B}

    \vspace{2pt}
    Claude 4.5 Sonnet modeli, özellikle kod yazma yeteneği, bağımsız çalışan yapay zeka ajanları (agents) oluşturma ve bilgisayar ekranını insan gibi kontrol edebilme (computer use) becerileriyle öne çıkmaktadır.
\end{tcolorbox}
\vspace{10pt}
% SORU 74
\needspace{5cm}
\vspace{8pt}
\noindent\textbf{\large\sffamily Soru 74.} \ Difüzyon (Diffusion) modellerinin temel çalışma prensibi aşağıdakilerden hangisidir?

\begin{enumerate}[label=\Alph*), leftmargin=25pt, topsep=4pt, itemsep=2pt]
        \item Metinlerdeki bir sonraki kelimeyi tahmin ederek ilerlemesi.
        \item İki sinir ağının birbiriyle yarışarak en gerçekçi sonucu bulması.
        \item Rastgele gürültü ile başlayıp bunu aşamalı olarak net bir görsele dönüştürmesi.
        \item Veriyi daha basit bir temsile sıkıştırıp sonra yeniden oluşturması.
\end{enumerate}

\vspace{4pt}
\begin{tcolorbox}[title=Cevap ve Açıklama, colback=gray!3!white, colframe=primarycolor!90!black, arc=2mm, fonttitle=\bfseries\sffamily, top=6pt, bottom=6pt]
    \textbf{Doğru Cevap: C}

    \vspace{2pt}
    Difüzyon (Diffusion) modelleri, rastgele gürültü içeren bir görselden yola çıkıp, ters difüzyon (denoising) adımlarıyla gürültüyü aşamalı olarak azaltarak net ve anlamlı görseller üretir.
\end{tcolorbox}
\vspace{10pt}
% SORU 75
\needspace{5cm}
\vspace{8pt}
\noindent\textbf{\large\sffamily Soru 75.} \ Gemini 2.5 Pro modelini rakiplerinden ayıran en önemli teknik özellik hangisidir?

\begin{enumerate}[label=\Alph*), leftmargin=25pt, topsep=4pt, itemsep=2pt]
        \item En yüksek parametre sayısına sahip olması.
        \item 1M token ile en geniş bağlam (kontekst) penceresini sunması.
        \item En düşük halüsinasyon oranına sahip olması.
        \item Açık kaynak kodlu tek lider model olması.
\end{enumerate}

\vspace{4pt}
\begin{tcolorbox}[title=Cevap ve Açıklama, colback=gray!3!white, colframe=primarycolor!90!black, arc=2mm, fonttitle=\bfseries\sffamily, top=6pt, bottom=6pt]
    \textbf{Doğru Cevap: B}

    \vspace{2pt}
    Gemini 2.5 Pro modelinin en belirgin teknik üstünlüğü, 1 milyon tokenlik devasa bağlam (context) penceresi sayesinde çok büyük veri dosyalarını tek seferde işleyebilmesidir.
\end{tcolorbox}
\vspace{10pt}
% SORU 76
\needspace{5cm}
\vspace{8pt}
\noindent\textbf{\large\sffamily Soru 76.} \ Transformer mimarisinin RNN (Tekrarlayan Sinir Ağları) gibi eski yöntemlere göre en büyük avantajı nedir?

\begin{enumerate}[label=\Alph*), leftmargin=25pt, topsep=4pt, itemsep=2pt]
        \item Eğitim için GPU donanımına ihtiyaç duymaması.
        \item Veri setindeki tüm öğeler arasındaki ilişkileri aynı anda (paralel) değerlendirebilmesi.
        \item Sadece sayısal veriler üzerinde uzmanlaşmış olması.
        \item Verileri sıralı bir şekilde tek tek okuması.
\end{enumerate}

\vspace{4pt}
\begin{tcolorbox}[title=Cevap ve Açıklama, colback=gray!3!white, colframe=primarycolor!90!black, arc=2mm, fonttitle=\bfseries\sffamily, top=6pt, bottom=6pt]
    \textbf{Doğru Cevap: B}

    \vspace{2pt}
    Transformer mimarisi, RNN'lerin aksine verileri sıralı değil, paralel olarak işler ve tüm dizi öğeleri arasındaki ilişkileri aynı anda değerlendirerek devasa bir eğitim avantajı sağlar.
\end{tcolorbox}
\vspace{10pt}
% SORU 77
\needspace{5cm}
\vspace{8pt}
\noindent\textbf{\large\sffamily Soru 77.} \ Yapay zeka eğitim sürecinde 'İnce Ayar' (Fine-tuning) aşamasının temel amacı nedir?

\begin{enumerate}[label=\Alph*), leftmargin=25pt, topsep=4pt, itemsep=2pt]
        \item Modeli tamamen sıfırdan eğiterek yeni bir dil öğretmek.
        \item Modelin belirli görevlerde daha başarılı olması ve çıktı güvenliğinin sağlanması.
        \item İnternetteki tüm verileri temizlemeden modele yüklemek.
        \item Modelin donanım hızını artırmak.
\end{enumerate}

\vspace{4pt}
\begin{tcolorbox}[title=Cevap ve Açıklama, colback=gray!3!white, colframe=primarycolor!90!black, arc=2mm, fonttitle=\bfseries\sffamily, top=6pt, bottom=6pt]
    \textbf{Doğru Cevap: B}

    \vspace{2pt}
    İnce Ayar (Fine-tuning), genel yeteneklere sahip bir ön eğitimli modeli (pre-trained), belirli bir sektörel alanda veya özel görevlerde daha yüksek performans ve güvenlik standartlarıyla çalıştırmak için ek verilerle eğitmektir.
\end{tcolorbox}
\vspace{10pt}
% SORU 78
\needspace{5cm}
\vspace{8pt}
\noindent\textbf{\large\sffamily Soru 78.} \ Aşağıdakilerden hangisi üretken yapay zekanın kullanımıyla ilgili belirtilen etik ve sosyal risklerden biridir?

\begin{enumerate}[label=\Alph*), leftmargin=25pt, topsep=4pt, itemsep=2pt]
        \item İstihdam endişeleri ve deepfake teknolojisi.
        \item Modellerin daha fazla dili desteklemesi.
        \item API kullanım maliyetlerinin düşmesi.
        \item Modellerin çok hızlı yanıt vermesi.
\end{enumerate}

\vspace{4pt}
\begin{tcolorbox}[title=Cevap ve Açıklama, colback=gray!3!white, colframe=primarycolor!90!black, arc=2mm, fonttitle=\bfseries\sffamily, top=6pt, bottom=6pt]
    \textbf{Doğru Cevap: A}

    \vspace{2pt}
    Üretken yapay zekanın en büyük sosyal ve etik riskleri arasında deepfake teknolojisi ile dezenformasyon yayılması ve çeşitli meslek dallarındaki istihdam endişeleri yer alır.
\end{tcolorbox}
\vspace{10pt}
% SORU 79
\needspace{5cm}
\vspace{8pt}
\noindent\textbf{\large\sffamily Soru 79.} \ 1950'lerden günümüze yapay zeka tarihindeki önemli dönüm noktalarından biri olan ve IBM tarafından geliştirilen 'Deep Blue'nun başarısı nedir?

\begin{enumerate}[label=\Alph*), leftmargin=25pt, topsep=4pt, itemsep=2pt]
        \item Dünya satranç şampiyonunu yenmesi.
        \item İlk chatbot olması.
        \item İnterneti kendi başına tarayabilen ilk yazılım olması.
        \item İnsan beynini tamamen taklit eden ilk sinir ağını kurması.
\end{enumerate}

\vspace{4pt}
\begin{tcolorbox}[title=Cevap ve Açıklama, colback=gray!3!white, colframe=primarycolor!90!black, arc=2mm, fonttitle=\bfseries\sffamily, top=6pt, bottom=6pt]
    \textbf{Doğru Cevap: A}

    \vspace{2pt}
    Yapay zeka tarihinin en önemli dönüm noktalarından biri, IBM tarafından geliştirilen 'Deep Blue' bilgisayarının 1997 yılında dünya satranç şampiyonu Garry Kasparov'u yenmesidir.
\end{tcolorbox}
\vspace{10pt}

\chapter{Genel Değerlendirme Sınav Soruları (Görseller)}
% SORU 80
\needspace{5cm}
\vspace{8pt}
\noindent\textbf{\large\sffamily Soru 80.} \ Prompt mühendisliği neden giderek daha önemli hale gelmektedir?

\begin{enumerate}[label=\Alph*), leftmargin=25pt, topsep=4pt, itemsep=2pt]
        \item Yapay zekanın doğru ve etkili sonuç üretmesi için
        \item Model eğitimi kolaylaştırmak için
        \item Sadece yazılı içerik için gerekli olduğu için
        \item Donanım üretimini artırmak için
\end{enumerate}

\vspace{4pt}
\begin{tcolorbox}[title=Cevap ve Açıklama, colback=gray!3!white, colframe=primarycolor!90!black, arc=2mm, fonttitle=\bfseries\sffamily, top=6pt, bottom=6pt]
    \textbf{Doğru Cevap: A}
    
    \vspace{2pt}
    Üretken yapay zeka modelleri (LLM'ler), girdilerin (prompt) formülasyonuna son derece duyarlıdır. Üstün kalitede, doğru ve işlevsel sonuçlar almak için doğru prompt tasarımı kritik öneme sahiptir.
\end{tcolorbox}
\vspace{10pt}
% SORU 81
\needspace{5cm}
\vspace{8pt}
\noindent\textbf{\large\sffamily Soru 81.} \ Yerelde çalışan bir yapay zeka modeli, genellikle daha yüksek gizlilik sağlar.

\begin{enumerate}[label=\Alph*), leftmargin=25pt, topsep=4pt, itemsep=2pt]
        \item false
        \item true
\end{enumerate}

\vspace{4pt}
\begin{tcolorbox}[title=Cevap ve Açıklama, colback=gray!3!white, colframe=primarycolor!90!black, arc=2mm, fonttitle=\bfseries\sffamily, top=6pt, bottom=6pt]
    \textbf{Doğru Cevap: B (true)}
    
    \vspace{2pt}
    Yerel çalışan modeller veriyi dış sunuculara veya buluta göndermediği için tam veri gizliliği sağlar.
\end{tcolorbox}
\vspace{10pt}
% SORU 82
\needspace{5cm}
\vspace{8pt}
\noindent\textbf{\large\sffamily Soru 82.} \ Üretken yapay zekada 'Gelecek Trendleri' konusu aşağıdakilerden hangisini kapsamaz?

\begin{enumerate}[label=\Alph*), leftmargin=25pt, topsep=4pt, itemsep=2pt]
        \item Model boyutlarının küçülmesi
        \item Yeni donanım gelişmeleri
        \item Yapay zeka problemlerini
        \item Yasal düzenlemeler
\end{enumerate}

\vspace{4pt}
\begin{tcolorbox}[title=Cevap ve Açıklama, colback=gray!3!white, colframe=primarycolor!90!black, arc=2mm, fonttitle=\bfseries\sffamily, top=6pt, bottom=6pt]
    \textbf{Doğru Cevap: C}
    
    \vspace{2pt}
    Yapay zeka problemleri güncel teknik sınırlılıklardır; gelecek yönelimli birer trend değillerdir.
\end{tcolorbox}
\vspace{10pt}
% SORU 83
\needspace{5cm}
\vspace{8pt}
\noindent\textbf{\large\sffamily Soru 83.} \ Üretken yapay zeka ile ilgili gelecek düzenlemelerin en çok hangi alanı kapsaması beklenir?

\begin{enumerate}[label=\Alph*), leftmargin=25pt, topsep=4pt, itemsep=2pt]
        \item Sosyal medya reklamcılığı
        \item Veri gizliliği ve telif hakları
        \item Arayüz düzenlemeleri
        \item Kullanıcı deneyimi
\end{enumerate}

\vspace{4pt}
\begin{tcolorbox}[title=Cevap ve Açıklama, colback=gray!3!white, colframe=primarycolor!90!black, arc=2mm, fonttitle=\bfseries\sffamily, top=6pt, bottom=6pt]
    \textbf{Doğru Cevap: B}
    
    \vspace{2pt}
    Modellerin devasa veriyle eğitilmesi fikri mülkiyet ve telif hakları tartışmalarını ön plana çıkarır; düzenlemeler bu alana odaklanır.
\end{tcolorbox}
\vspace{10pt}
% SORU 84
\needspace{5cm}
\vspace{8pt}
\noindent\textbf{\large\sffamily Soru 84.} \ Üretken yapay zekanın en temel kullanım alanlarından biri aşağıdakilerden hangisidir?

\begin{enumerate}[label=\Alph*), leftmargin=25pt, topsep=4pt, itemsep=2pt]
        \item İçerik üretimi
        \item Donanım üretimi
        \item Veri depolama
        \item Sunucu yönetimi
\end{enumerate}

\vspace{4pt}
\begin{tcolorbox}[title=Cevap ve Açıklama, colback=gray!3!white, colframe=primarycolor!90!black, arc=2mm, fonttitle=\bfseries\sffamily, top=6pt, bottom=6pt]
    \textbf{Doğru Cevap: A}
    
    \vspace{2pt}
    Üretken yapay zekanın temel varlık amacı yeni ve özgün metin, resim, ses vb. içerikler üretmektir.
\end{tcolorbox}
\vspace{10pt}
% SORU 85
\needspace{5cm}
\vspace{8pt}
\noindent\textbf{\large\sffamily Soru 85.} \ Gelecekte üretken yapay zekaya ilişkin regülasyonlar neden önem kazanacaktır?

\begin{enumerate}[label=\Alph*), leftmargin=25pt, topsep=4pt, itemsep=2pt]
        \item Güvenlik, gizlilik ve etik sorunları kontrol altına almak için
        \item Yapay zekanın hızını düşürmek için
        \item Donanım üretimini artırmak için
        \item Modelin performansını artırmak için
\end{enumerate}

\vspace{4pt}
\begin{tcolorbox}[title=Cevap ve Açıklama, colback=gray!3!white, colframe=primarycolor!90!black, arc=2mm, fonttitle=\bfseries\sffamily, top=6pt, bottom=6pt]
    \textbf{Doğru Cevap: A}
    
    \vspace{2pt}
    Yapay zekanın kontrolsüz gelişimini, etik dışı kullanımları ve güvenlik risklerini yasal sınırlarla yönetmek regülasyonların temel amacıdır.
\end{tcolorbox}
\vspace{10pt}
% SORU 86
\needspace{5cm}
\vspace{8pt}
\noindent\textbf{\large\sffamily Soru 86.} \ Üretken yapay zeka modellerinde 'zero-shot learning' özelliği ne sağlar?

\begin{enumerate}[label=\Alph*), leftmargin=25pt, topsep=4pt, itemsep=2pt]
        \item Modelin hiç eğitim almadan yeni görevleri çözebilmesi
        \item Modelin verisiz çalışması
        \item Modelin hızlanması
        \item Modelin hata yapması
\end{enumerate}

\vspace{4pt}
\begin{tcolorbox}[title=Cevap ve Açıklama, colback=gray!3!white, colframe=primarycolor!90!black, arc=2mm, fonttitle=\bfseries\sffamily, top=6pt, bottom=6pt]
    \textbf{Doğru Cevap: A}
    
    \vspace{2pt}
    Zero-shot, modelin daha önce özel eğitim görmediği yepyeni görevleri ön eğitim becerileriyle çözmesidir.
\end{tcolorbox}
\vspace{10pt}
% SORU 87
\needspace{5cm}
\vspace{8pt}
\noindent\textbf{\large\sffamily Soru 87.} \ Bir 'Prompt' optimize edilirken aşağıdakilerden hangisine dikkat edilmelidir?

\begin{enumerate}[label=\Alph*), leftmargin=25pt, topsep=4pt, itemsep=2pt]
        \item Rastgele sözcüklerle yazılmasına
        \item Sadece İngilizce yazılmasına
        \item Çok uzun ve karmaşık olmasına
        \item Kısa ve öz olmasına
\end{enumerate}

\vspace{4pt}
\begin{tcolorbox}[title=Cevap ve Açıklama, colback=gray!3!white, colframe=primarycolor!90!black, arc=2mm, fonttitle=\bfseries\sffamily, top=6pt, bottom=6pt]
    \textbf{Doğru Cevap: D}
    
    \vspace{2pt}
    Aşırı karmaşıklıktan uzak, açık, net, kısa ve öz girdiler yapay zekadan en yüksek performansı almamızı sağlar.
\end{tcolorbox}
\vspace{10pt}
% SORU 88
\needspace{5cm}
\vspace{8pt}
\noindent\textbf{\large\sffamily Soru 88.} \ Vektör veritabanları doğrudan metni saklar, vektöre dönüştürmez.

\begin{enumerate}[label=\Alph*), leftmargin=25pt, topsep=4pt, itemsep=2pt]
        \item false
        \item true
\end{enumerate}

\vspace{4pt}
\begin{tcolorbox}[title=Cevap ve Açıklama, colback=gray!3!white, colframe=primarycolor!90!black, arc=2mm, fonttitle=\bfseries\sffamily, top=6pt, bottom=6pt]
    \textbf{Doğru Cevap: A (false)}
    
    \vspace{2pt}
    Vektör veritabanları metni doğrudan saklamaz, öncelikle yüksek boyutlu sayısal vektörlere (embedding) dönüştürüp saklar.
\end{tcolorbox}
\vspace{10pt}
% SORU 89
\needspace{5cm}
\vspace{8pt}
\noindent\textbf{\large\sffamily Soru 89.} \ Prompt mühendisliğinde aşağıdakilerden hangisi önemlidir?

\begin{enumerate}[label=\Alph*), leftmargin=25pt, topsep=4pt, itemsep=2pt]
        \item Komutu karmaşık hale getirmek
        \item Komutu rastgele oluşturmak
        \item Çok fazla onay parametresi belirtmek
        \item Açık ve net bir komut vermek
\end{enumerate}

\vspace{4pt}
\begin{tcolorbox}[title=Cevap ve Açıklama, colback=gray!3!white, colframe=primarycolor!90!black, arc=2mm, fonttitle=\bfseries\sffamily, top=6pt, bottom=6pt]
    \textbf{Doğru Cevap: D}
    
    \vspace{2pt}
    Prompt mühendisliğinin altın kuralı doğrudan, açık ve net komutlar vermektir.
\end{tcolorbox}
\vspace{10pt}
% SORU 90
\needspace{5cm}
\vspace{8pt}
\noindent\textbf{\large\sffamily Soru 90.} \ Prompt mühendisliği yalnızca eğitim aşamasında kullanılır.

\begin{enumerate}[label=\Alph*), leftmargin=25pt, topsep=4pt, itemsep=2pt]
        \item false
        \item true
\end{enumerate}

\vspace{4pt}
\begin{tcolorbox}[title=Cevap ve Açıklama, colback=gray!3!white, colframe=primarycolor!90!black, arc=2mm, fonttitle=\bfseries\sffamily, top=6pt, bottom=6pt]
    \textbf{Doğru Cevap: A (false)}
    
    \vspace{2pt}
    Prompt mühendisliği model eğitiminde değil, eğitilmiş modelden çıktı alırken (inference) kullanılır.
\end{tcolorbox}
\vspace{10pt}
% SORU 91
\needspace{5cm}
\vspace{8pt}
\noindent\textbf{\large\sffamily Soru 91.} \ Üretken yapay zeka ile yapılan içerik üretiminde telif hakkı problemi neden önemlidir?

\begin{enumerate}[label=\Alph*), leftmargin=25pt, topsep=4pt, itemsep=2pt]
        \item İçeriklerin hızlı üretilmesi için
        \item İçeriklerin küçültülmesi için
        \item Donanım tasarımı için
        \item İçeriklerin yasal olarak korunması gerektiği için
\end{enumerate}

\vspace{4pt}
\begin{tcolorbox}[title=Cevap ve Açıklama, colback=gray!3!white, colframe=primarycolor!90!black, arc=2mm, fonttitle=\bfseries\sffamily, top=6pt, bottom=6pt]
    \textbf{Doğru Cevap: D}
    
    \vspace{2pt}
    Yapay zeka ürünlerinin yasal sahipliği ve yasal koruma altına alınması, ticari hak ihlallerinin önlenmesi açısından telif hakları en önemli konudur.
\end{tcolorbox}
\vspace{10pt}
% SORU 92
\needspace{5cm}
\vspace{8pt}
\noindent\textbf{\large\sffamily Soru 92.} \ Retrieval Augmented Generation (RAG) modelinde veri nereye bakılarak üretilir?

\begin{enumerate}[label=\Alph*), leftmargin=25pt, topsep=4pt, itemsep=2pt]
        \item Veritabanı indekslerine
        \item Dış bilgi kaynağından getirilen içeriklere
        \item Sabit kod bloklarına
        \item Model parametrelerine
\end{enumerate}

\vspace{4pt}
\begin{tcolorbox}[title=Cevap ve Açıklama, colback=gray!3!white, colframe=primarycolor!90!black, arc=2mm, fonttitle=\bfseries\sffamily, top=6pt, bottom=6pt]
    \textbf{Doğru Cevap: B}
    
    \vspace{2pt}
    RAG sistemleri, dış bilgi kaynaklarından dinamik veri çekerek LLM'lerin güncel kalmasını ve bağlama dayalı doğru üretim yapmasını sağlar.
\end{tcolorbox}
\vspace{10pt}
% SORU 93
\needspace{5cm}
\vspace{8pt}
\noindent\textbf{\large\sffamily Soru 93.} \ Bir 'Prompt' aşağıdakilerden hangisini ifade eder?

\begin{enumerate}[label=\Alph*), leftmargin=25pt, topsep=4pt, itemsep=2pt]
        \item Yapay zekaya verilen yönerge veya girdiyi
        \item Yapay zeka modelinin mimarisini
        \item Sonuç verilerini düzenleme yöntemini
        \item Modelin eğitim verilerini
\end{enumerate}

\vspace{4pt}
\begin{tcolorbox}[title=Cevap ve Açıklama, colback=gray!3!white, colframe=primarycolor!90!black, arc=2mm, fonttitle=\bfseries\sffamily, top=6pt, bottom=6pt]
    \textbf{Doğru Cevap: A}
    
    \vspace{2pt}
    Prompt, modele yöneltilen her türlü yazılı talimat, soru veya girdidir.
\end{tcolorbox}
\vspace{10pt}
% SORU 94
\needspace{5cm}
\vspace{8pt}
\noindent\textbf{\large\sffamily Soru 94.} \ Üretken yapay zeka teknolojilerinin geleceğinde düzenlemelerin artması beklenmektedir.

\begin{enumerate}[label=\Alph*), leftmargin=25pt, topsep=4pt, itemsep=2pt]
        \item false
        \item true
\end{enumerate}

\vspace{4pt}
\begin{tcolorbox}[title=Cevap ve Açıklama, colback=gray!3!white, colframe=primarycolor!90!black, arc=2mm, fonttitle=\bfseries\sffamily, top=6pt, bottom=6pt]
    \textbf{Doğru Cevap: B (true)}
    
    \vspace{2pt}
    Yapay zeka güvenliği ve dezenformasyonun önlenmesi gibi amaçlarla yasal denetimlerin artacağı kesin olarak öngörülmektedir.
\end{tcolorbox}
\vspace{10pt}
% SORU 95
\needspace{5cm}
\vspace{8pt}
\noindent\textbf{\large\sffamily Soru 95.} \ Dil modellerinde 'token' kavramı neyi ifade eder?

\begin{enumerate}[label=\Alph*), leftmargin=25pt, topsep=4pt, itemsep=2pt]
        \item Veri tabanındaki dosyayı
        \item Bir kelimenin anlamını
        \item Bir cümlenin tamamını
        \item Bir kelimenin ya da parçanın sayısal karşılığını
\end{enumerate}

\vspace{4pt}
\begin{tcolorbox}[title=Cevap ve Açıklama, colback=gray!3!white, colframe=primarycolor!90!black, arc=2mm, fonttitle=\bfseries\sffamily, top=6pt, bottom=6pt]
    \textbf{Doğru Cevap: D}
    
    \vspace{2pt}
    Dil modelleri metinleri doğrudan okuyamaz. Öncelikle her kelime veya hece token adı verilen ve sayısal karşılığı (ID) olan parçalara bölünür.
\end{tcolorbox}
\vspace{10pt}
% SORU 96
\needspace{5cm}
\vspace{8pt}
\noindent\textbf{\large\sffamily Soru 96.} \ Üretken yapay zekanın gelecekte en çok etkileyebileceği alanlardan biri aşağıdakilerden hangisidir?

\begin{enumerate}[label=\Alph*), leftmargin=25pt, topsep=4pt, itemsep=2pt]
        \item Donanım üretimi
        \item Blockchain güvenliği
        \item Otonom içerik üretimi
        \item Biyolojik veri şifreleme
\end{enumerate}

\vspace{4pt}
\begin{tcolorbox}[title=Cevap ve Açıklama, colback=gray!3!white, colframe=primarycolor!90!black, arc=2mm, fonttitle=\bfseries\sffamily, top=6pt, bottom=6pt]
    \textbf{Doğru Cevap: C}
    
    \vspace{2pt}
    Metin, görsel, ses ve video üretme yeteneği sayesinde otonom içerik üretimi süreçleri en çok etkilenecek sektördür.
\end{tcolorbox}
\vspace{10pt}
% SORU 97
\needspace{5cm}
\vspace{8pt}
\noindent\textbf{\large\sffamily Soru 97.} \ Yapay zeka destekli ofis araçları ile aşağıdakilerden hangisi yapılabilir?

\begin{enumerate}[label=\Alph*), leftmargin=25pt, topsep=4pt, itemsep=2pt]
        \item Donanım üretimi
        \item Ağ yönetimi
        \item Veri şifreleme
        \item Metin ve sunum hazırlama
\end{enumerate}

\vspace{4pt}
\begin{tcolorbox}[title=Cevap ve Açıklama, colback=gray!3!white, colframe=primarycolor!90!black, arc=2mm, fonttitle=\bfseries\sffamily, top=6pt, bottom=6pt]
    \textbf{Doğru Cevap: D}
    
    \vspace{2pt}
    Modern yapay zeka ofis araçları (Copilot vb.) metin yazımı, veri analizi ve otomatik sunum hazırlama gibi işlerde etkin rol oynar.
\end{tcolorbox}
\vspace{10pt}
% SORU 98
\needspace{5cm}
\vspace{8pt}
\noindent\textbf{\large\sffamily Soru 98.} \ Üretken yapay zekanın dezavantajlarından biri aşağıdakilerden hangisidir?

\begin{enumerate}[label=\Alph*), leftmargin=25pt, topsep=4pt, itemsep=2pt]
        \item Çok düşük veri kullanımı
        \item Yüksek donanım dayanıklılığı
        \item Yanlış veya yanıltıcı bilgi üretimi
        \item Daha az esneklik
\end{enumerate}

\vspace{4pt}
\begin{tcolorbox}[title=Cevap ve Açıklama, colback=gray!3!white, colframe=primarycolor!90!black, arc=2mm, fonttitle=\bfseries\sffamily, top=6pt, bottom=6pt]
    \textbf{Doğru Cevap: C}
    
    \vspace{2pt}
    Yapay zeka modellerinin tamamen uydurma ve yanlış bilgileri son derece ikna edici biçimde üretmesine halüsinasyon denir; en büyük dezavantajlardan biridir.
\end{tcolorbox}
\vspace{10pt}
% SORU 99
\needspace{5cm}
\vspace{8pt}
\noindent\textbf{\large\sffamily Soru 99.} \ Yapay zeka destekli ofis araçları, sadece yazı yazmak için kullanılabilir.

\begin{enumerate}[label=\Alph*), leftmargin=25pt, topsep=4pt, itemsep=2pt]
        \item false
        \item true
\end{enumerate}

\vspace{4pt}
\begin{tcolorbox}[title=Cevap ve Açıklama, colback=gray!3!white, colframe=primarycolor!90!black, arc=2mm, fonttitle=\bfseries\sffamily, top=6pt, bottom=6pt]
    \textbf{Doğru Cevap: A (false)}
    
    \vspace{2pt}
    Yapay zeka destekli ofis araçları yazı dışında veri analizi yapma, sunum hazırlama, grafik çizme ve özet çıkarma gibi birçok karmaşık görevi yapabilir.
\end{tcolorbox}
\vspace{10pt}

\end{document}
